\section{From a finite model to a deployable policy}\label{app:extensions}

The next results change the granularity of the admitted evidence. A source program can describe many contracts compactly. A receipt from performing one job can constrain which jobs are lawful later. A statistical record can carry information in its joint distribution that its individual fields lose. In each case the question is which distinctions a sound decision procedure must retain.

\subsection{Product-stage deadlines}

Let \(n\ge1\). The installed service has modes \(m=(m_1,\ldots,m_n)\) in a nonempty finite Cartesian class. Each query is a pure, lawful, deterministic read of one stage at a fixed positive rational price. Its answer restricts that stage alone. Write \(U=\prod_i U_i\) for a nonempty observed cell. There are \(n\) mandatory jobs in fixed order, an observed stage index, and no useful shortcut. For mode bit \(x_i(m_i)\) the fast atom survives the representation freeze exactly when the bit is zero. Its tariff is \(a_i\); the always-lawful slow atom costs \(a_i+w_i\), with \(a_i,w_i\) positive integers. Every supported job outcome actually completes the stage and advances. Concrete implementation anchors cover the complete callable action interface, conditional enablement, tariffs, extensional touch, positive support, public records and physical job events. The same public command must implement a chosen move across every hidden checkpoint in the current belief. Knowing the conditional interface does not reveal the installed mode. Apart from priced reads, earlier history does not constrain later residual modes.

\begin{theorem}[T30.S. Product-stage deadline]\label{ext:t30s}

Define \(v(m)=\sum_i(a_i+w_ix_i(m_i))\), the least worst-supported tariff for serving the installed system. Over every nonempty Cartesian read cell \(U\), one public policy achieves worst tariff \(\max_Uv\) and no common policy can improve it. Thus deadline \(B\ge0\) is certified exactly if \(\max_Uv\le B\), refuted exactly if \(\min_Uv>B\), and otherwise has witnesses to both verdicts.

A pair of modes at stage \(i\) crosses if two installations differing only by that pair fall on opposite sides of \(B\). The minimum cost of a uniform static certificate on \(U\) is the sum, over stages, of the cheapest sets of local queries separating all local crossing pairs. An unseparated pair under the full menu witnesses impossibility even for adaptive acquisition. With at most \(d\) modes and \(q\) queries per stage and \(b=\max(0,B-\sum_i a_i)\), a direct implementation uses
\[
O\!\left(n^2d(b+2)+nd^2(q+1)(b+2)+n2^qd^2(q+1)\right)
\]
unit operations, with integer sums capped at \(b+1\) in the decision calculation. In the gateway family with \(n\) ancillary jobs below, the target and full-mode costs are \(1\) and \(n+1\). The result concerns pure static read prices and worst-supported service. Zero probability of deadline failure is equivalent to sure success separately for each installation when its whole support has positive probability.

\end{theorem}

\begin{proof}

\textbf{Premises.}

A stage has nonempty finite mode set \(D_i\), fast tariff \(a_i>0\), extra slow tariff \(w_i>0\) and mode bit \(x_i\). The frozen contract admits fast exactly when \(x_i=0\). Slow is always legal. Both complete the job on their whole positive support and advance. Other acts cannot complete or skip this mandatory service. Costs are finite integers. The stage index and prior legitimate query answers are observable. Queries are source-grounded readbacks about one stage, with disjoint namespaces and fixed positive rational prices. A given answer cell is \(U=\prod_i U_i\) because query restrictions are stage-local, with nonempty \(U_i\) checked.

\textbf{A. Transport and the common policy.}

At stage \(i\) of contract \(m\), any successful policy must perform the local job at cost at least \(a_i+w_i x_i(m_i)\). Taking fast exactly when it survives the freeze achieves this price at every stage. Induction on the required job chain proves the per-contract optimum \(v(m)=\sum_i(a_i+w_i x_i(m_i))\). Positive support includes every possible successor declared by the contract. In particular a fast action with one R-changing successor disappears as an entire atom.

The concrete deployment checkpoint map starts at phase 1. At each physical checkpoint the anchored row fixes the complete exposed action classes, R-survival, tariff, positive support, next-phase projection and actual job event. The locally required actions and nonempty supports exist. Each required canonical action has a common public callable lift across all hidden states and residual modes where that action is used. The lift is a function of publicly available action/history information and the audited cell. Independently checked row identities hold conditionally for each permitted mode, without releasing the installed mode as an unpaid observation. Induct on the actual concrete steps. Every branch of the scheduled controller executes exactly one projected job in the next required phase, incurs the stated tariff, and progresses to done. No other concrete lawful execution can improve the lower bound, since the entire exposed job menu is represented. This derives the deployment optimum from guard, tariff, support and job-event checks that do not require a deadline verdict. The physical audit must justify that its inventory includes every relevant supported edge. A verified computation on a declared inventory supplies no missing physical evidence.

For a nonempty cell put \(h_i=\max_{m_i\in U_i}x_i(m_i)\). One observable schedule uses fast at \(i\) if \(h_i=0\) and slow otherwise. It is lawful throughout the cell and performs every job. Its worst cost is \(\sum_i(a_i+w_i h_i)\). Choose each local maximizing mode and take their product tuple. That tuple exists by independence of the compatibility domains and its known-mode optimum is the same sum. Every common policy has worst cost at least this much. The equality with \(\max_U v\) is therefore constructive. It does not infer a common policy simply from separate existence. In a finite stochastic interpretation with the same coherent positive support, all-path success is equivalent to zero failure probability. If a supported path fails, that path has positive mass under a law in the declared support class. This proves a zero-risk family statement without estimating positive risk from support.

\textbf{B. Exact target set and witnesses.}

Get \(L=\sum_i(a_i+w_i\min_{U_i}x_i)\) and \(H=\sum_i(a_i+w_i\max_{U_i}x_i)\), after testing nonemptiness. The extrema are attained by actual product tuples. If \(H\le B\) the target equals 1 throughout and the schedule above is the shared positive witness. If \(L>B\) it equals 0 throughout, because even fully informed controls miss the deadline. Otherwise \(L\le B<H\) and the extrema give an explicit target-distinct pair. Equality at the deadline passes. An empty answer cell is inconsistent. Refusal fails the all-jobs requirement.

\textbf{C. Minimum static audit and a full-menu obstruction.}

A pair of local modes \(u,u'\) for stage \(i\) crosses if the two deadline tests disagree when every other stage has exactly the same fixed choices. Its edge is the set of stage-\(i\) queries with unequal answers. Necessity of hitting this edge follows from the corresponding actual compatible pair: queries elsewhere have identical answers. For sufficiency take any distinct-target completions with the same selected answers and pass between them by replacing modes one coordinate at a time. Every intermediate tuple belongs to the product. Each selected readback is local, so its value stays constant along the hybrid path. At one step the target changes, exposing a missed local edge. Thus separating every local crossing is exactly uniform identification.

Disjoint namespaces and additive prices imply a lower bound given by the sum of local optimum hitting costs. Take one local minimum for each stage and their union attains the bound. All leaves of this static audit are homogeneous and the positive leaves are Cartesian cells that admit the common schedule. If a local edge is empty, its witnessing tuples agree on the full legal menu. Inducting on acquired answers shows that every adaptive strategy has the same history on these two modes for each fixed random tape. It cannot certify both correctly with zero error. If all local edges are nonempty, buying all allowed readbacks is feasible and an exact optimizer finds the least static cost. Audits with history-dependent availability or prices are outside this optimization.

\textbf{D. Numeric representation and complexity.}

Set \(A=\sum_i a_i\). If \(B<A\) the target is constant false on a nonempty cell. Otherwise let \(b=B-A\). For each excluded stage \(i\), compute outside excess sums using \(x\oplus y=\min(b+1,x+y)\). Nonnegative excesses make capping commute with further additions. Start with sum zero and add each remaining stage's declared possible increments. Keep one predecessor and actual mode choice per attained capped sum per layer. Induction on layers proves that every stored value has a compatible preimage and every real attainable sum is represented by its cap. Testing \(s+w_i x_i(u)\le b\) against the other local mode's test supplies exactly the crossings, together with compatible outside witnesses. The overflow value cannot falsely pass a test.

There are at most \(b+2\) sums per layer. Independently discarding each stage costs \(O(n^2d(b+2))\) updates for maximum local domain size \(d\). Pair checking and edge formation take at most \(O(nd^2(q+1)(b+2))\) for maximum local menu \(q\). Enumeration of local query subsets costs \(O(n2^q d^2(q+1))\). Working tables with backpointers for one exclusion use \(O(n(b+2))\) space, in addition to inputs and retained edges/witnesses. Backtracing a chosen witness is linear in stages. Exact rational acquisition requires standard exact bit arithmetic for tariffs and prices. These counts concern the already verified stage input and do not account for the independent evidentiary work of verifying the deployment anchor. There is no explicit global completion catalog. The bound is pseudo-polynomial in the numeric deadline and exponential in local menu size. Fixed \(d,q\) and polynomial numeric \(b\) yield a polynomial structural algorithm even for an exponentially large product. This comparison is with explicit materialization, with no universal runtime comparison to other structured solvers.

\textbf{E. Non-vacuous PC example.}

Use two worlds with distinct sealed targets. Every pre-stage history \(p_i\) is compatible with both worlds. Terminal evidence records \(\tau_0,\tau_1\) are individually compatible with their indexed worlds. Also use paired legitimate job-certificate histories \(b_{i,0},b_{i,1}\). Make the pair one extensional certificate block exactly when stale, while progress and final evidence records are singleton blocks. The source action fast performs job \(i\) and refreshes its stale block by splitting it when needed. Slow performs the same job without changing the relation. In either case the emitted admitted progress advances to the next \(p_i\), or finally a world-specific \(\tau_u\). All extensional changes are derived over the complete source atom and its support. Authority is fixed. In the original PC closure sense the initial represented fiber remains open and the terminal checkpoints are individually closed. If terminal reporting is required, a final visible result can disclose the sealed value without affecting the path cost. Each mode bit genuinely changes whether the fast service route survives the R freeze. This is actual work with an independently meaningful representation state, rather than a collection of inspection-only dummy fields.

For \(n\) ancillary jobs plus one gateway, take all base tariffs 1 and ancillary extra tariffs 1. Set gateway extra tariff \(n+1\) and deadline \(2n+1\). A fresh gateway passes even when every ancillary mode is stale, since \((n+1)+n=2n+1\). A stale gateway fails even when all others are fresh, since \((n+1)+(n+1)>2n+1\). Each local mode has an independent unit-price lawful readback. A single gateway read suffices and zero reads cannot resolve both target values. Full identity requires every bit, since each omitted bit can change independently and genuinely alters a required job's R-frozen fast service. The cost ratio is \(n+1\) in the specified family. In the positive cell a single feasible schedule serves all actual jobs.

The same construction supplies the deadline verdict, attaining installations, least-cost static read set and common service policy. The price lower bound follows from actual crossing pairs. The concrete service guarantee follows from the independently checked physical rows. Local capped sums establish the computational bound without a catalog of all installations.

The Cartesian and stage-local hypotheses are essential to this argument. Theorem~\ref{ext:t30cg} allows correlated installations in a different online protocol. The construction in Theorem~\ref{adv:ug3} gives statistical and strategic guarantees for its own checked finalizer and observation menu. Neither result gives the Cartesian price formula for arbitrary interacting services.

\end{proof}

\subsection{M01. Finite source interpretation and observable simulation}

\label{ext:m01}

The input to the first construction is an interpreted finite source, including its rules for which successors have positive weight. This allows replay of a witness down to the installed source slots. A second construction asks how much of this structure must be preserved to transport a policy.

\begin{theorem}[M01.1. finite typed source extraction and exact certificates]\label{ext:m011}

\textbf{Fragment.} Fix finite explicitly represented types \(X,U,Z,O,V\), a finite registry \(A\) of distinct whole source actions and an extensional action permission type \(2^A\). Fix a sealed world target \(A_\Pi:X\to D\) and a fixed entry locator. A checkpoint is a vector \((z,h,p,\lambda)\in S=Z\times U\times \operatorname{Eq}(U)\times 2^A\), with equality on the last two coordinates given by their extensional tables. The expansion of these types is a permitted computation cost. A finite tuple \(\theta\) in a product of finite interpreted slot domains is fixed for the entire contract. It can fill declared source field slots and environment parameters; its correlations are expressed by governing constraints. Typed, total straight-line circuits over these finite types give the entry \(s_0(\theta)\), compatibility \(x\sim_\theta h\), public controller observation \(\omega_\theta(s)\), and enabled-action and postcondition semantics as follows.

\begin{itemize}

\item
  A single registered action \(a\) is admitted exactly when its source guard (a total circuit with no permission-register argument) holds and \(a\in\lambda\). A declared finite collection of guarded environment alternatives gives whole postcheckpoint vectors. Exactly alternatives with true guard and strictly positive weight contribute to support. Weights are exact declared nonnegative rationals. A total outcome circuit \(Y_\theta(s,a,u)\in V\) gives the unique visible edge label for an ordered triple of checkpoints and action. Alternatively a source that retains an outcome in its raw state can use that component to evaluate this circuit. A finite nonnegative rational \(c(a)\) is fixed per source action. Controller choices use only the full history of observations, chosen action identities and visible outcomes. Controller observation and certificate equality are separate types.
\item
  A permitted machine has nonempty support at each admitted occurrence reachable from the entry and has a nonempty represented fiber at each reachable checkpoint. Reachability uses the unfrozen full positive support. The fiber at \(s\) is \(\{x:\exists h'\in [h]_p,\ x\sim_\theta h'\}\), and closure requires constancy of the sealed \(A_\Pi\) there. From the source support, the mask of an admitted occurrence is the union of extensional changes in \(h,p,\lambda\) over its supported successors. Freezing \(F\subseteq\{E,R,A\}\) removes a whole occurrence when its mask intersects \(F\). Kept occurrences retain every supported successor. The target \(\Phi\) is a specified total circuit over complete canonical kernel fields and, if requested, the exact eight-coordinate optimal proper closure signature, with \(+\infty\) represented by a separate flag.
\item
  Governing evidence \(\Omega\) is a Boolean circuit over the slots and the canonical fields, with allowed atoms including existential support of a particular finite trusted trace (truncated or complete as declared) and exact finite readbacks from named licensed source instruments. Its reachability or trace atoms inspect only source-generated entries and edges. A specified instrument returns a source-grounded finite value by a total circuit and does not change the plant, slots or target. Instrument outputs do not become controller observations unless explicitly in the controller interface. Requirements outside these finite typed schemas are outside the fragment. The actual system inclusion claim is conditional on the actual governing machine having this source semantics for some permitted \(\theta\) and satisfying \(\Omega\).
\end{itemize}

For every such source, there is a terminating typed-to-Boolean compilation of a reusable canonical kernel together with witness and refutation checking for the machine-validity and evidence predicate. For every assignment to its input slots, the kernel fields are uniquely determined by the source equations. Satisfying assignments to validity plus evidence project exactly to compatible interpretations, allowing duplicate representations that induce the same contract. Kernel evaluation preserves full support, availability, the observation boundary, the extensional resource masks and the eight finite-tree optima including \(+\infty\). Consequently the following tasks have total certificate-producing procedures for the declared encoding.

\begin{enumerate}
\def\labelenumi{\arabic{enumi}.}

\item
  Find a compatible slot witness or certify inconsistency. Certify the possible values (jointly, for a tuple) of the finite target \(\Phi\), and certify singleton identification only with a compatible witness and a checked no-opposite-value certificate. The same applies to a consistent observed readback cell. Finite minima/maxima here are attained.
\item
  Replay a paired ambiguity witness including the actual source assignments, exact targets and any readback equalities. Given a specified complete list of deterministic pure readbacks, either certify that their joint answers separate every opposite-target pair or exhibit two compatible opposite-target assignments with equal answers to the entire list. In the latter case no policy using solely these readbacks can identify uniformly with certainty, including adaptive selection.
\item
  For a fixed compatible kernel, a nonnegative integer budget \(b\), and any subset of the fixed-kernel admitted occurrences marked kept, construct a finite monotone Boolean circuit in the keep bits for proper closure by budget \(b\). In this clause costs are nonnegative integers. If the number of physical checkpoints is \(n\), using \(2^n-1\) as a depth bound suffices. An occurrence that is unavailable at any possible current checkpoint cannot serve as a common action. The inclusion-minimal satisfying keep sets are exactly the inclusion-minimal proper policy supports at budget \(b\). Replacing keep bits by the corresponding whole-atom resource-freeze bits specializes the circuit to each fixed \(F\).
\item
  All the preceding certificate tasks may use explicit finite expansion with local equations and a finite Boolean proof of absence. There is no bound polynomial in succinct source length in this claim. A solver timeout or a checkable satisfying relaxation without a compatible concrete source assignment is insufficient for exact identification, refutation or a concrete disagreement witness.
\end{enumerate}

\end{theorem}

\begin{theorem}[M01.2. directed observable belief cover]\label{ext:m012}

Fix two D-35 finite typed kernels with supported successors, nonnegative finite action costs, and nonempty observable entry beliefs. At each nonempty belief \(B\), define \(L(B)\) to be its whole source actions legal at every \(s\in B\). For \(a\in L(B)\), let \(U(B,a)\) be the union of canonical occurrence masks across \(B\). A child is the full nonempty successor belief \(B_{a,\ell}\) indexed by a feasible visible pair \(\ell=(Y,\omega_{\mathrm{next}})\). Write \(J(B,a)\) for the set of these labels. Declare \(D(B)\) to be its stopping menu: a single dummy label iff every state is closed for plain closure, or the set of permitted common report labels under an explicitly declared common-label stopping rule. In the second system use primes. The entry beliefs are reached by controllers that know their own histories, with no hidden state access.

A directed relation \(R\) on pairs of reachable unfrozen beliefs, with source first and receiver second, is a cover if for all \(B\,R\,B'\):

\begin{itemize}

\item
  \(D(B)\subseteq D'(B')\).
\item
  For every \(a\in L(B)\), there exists a \textbf{single} \(a'\in L'(B')\), with \(c'(a')\le c(a)\) and \(U'(B',a')\subseteq U(B,a)\), and there exists a function \(\sigma:J'(B',a')\to J(B,a)\) such that \(B_{a,\sigma(\ell')}\,R\,B'_{a',\ell'}\) for every \(\ell'\) in the domain.
\end{itemize}

The action and function may depend on the pair of belief nodes and the chosen source action. They do not depend on an unobserved physical state or successor. Commands remain actual source atoms in each system. No injectivity or global relabeling of menus or outcomes is required.

If the entry pair belongs to a cover, then for every \(F\subseteq\{E,R,A\}\) and every finite proper source policy after freeze \(F\), there is a finite proper receiving policy after the same freeze, with no greater depth or worst-support cost. The receiving controller can implement it from its own observable history. Thus \(\kappa_{\rm receiver}^{\neg F}\le\kappa_{\rm source}^{\neg F}\) in the extended reals. A directed cover in both directions gives equality of all eight coordinates, equality of defined \(R\)-penalties, and agreement on baseline unclosability. If a symmetric correspondence in the accepted T03 sense holds at the original entries, the lifted covered-belief pairs satisfy the directed requirements in each direction. The converse fails for finite typed contracts, and equality of signatures alone does not imply covers. On the explicit unfrozen belief graphs the greatest directed cover is computable and checkable by finite relation elimination in polynomial time in their expanded sizes. If a directed cover additionally admits a surjective \(\sigma\) and equality of move costs for each matching choice, it transports each proper source policy at exactly its worst-path cost; the respective two-way condition gives equality of attainable proper worst-path cost sets.

\textbf{Conditions and boundary cases.} Closure and report moves are separately covered. Applying T03 requires preservation of the actionable stopping and reporting menus and a common callable command. Equal closed hidden states can have different prerequisites for announcing a result. Pointwise menu correspondence cannot transport one common action.

\end{theorem}

\begin{theorem}[M01.3. interpretation boundaries]\label{ext:m013}

A deterministic typed source program with a deployed receiver and a lawful readback has two compatible slot assignments with the same fixed source program and identical complete traces in its permitted evaluation setting, and with opposite values of one fixed binary deployment failure target. The evaluating setting does not select the installed receiver record. A read-only measurement of this record separates the pair, and fixing it within that two-assignment evidence class identifies the target. Both assignments use the same independently fixed criterion for successful physical receipt.

A terminating exact compiler deciding the existence of a finite proper closer for all valid unbounded-memory, deterministic small-step sources would decide machine halting, even when there is a unique interpretation and the controller has a single possible tick action while running. This claim concerns extraction/evaluation of stronger source languages. The finite compiler is not asserted to infer that a physical installation matches a chosen source language or slot schema.

\end{theorem}

\begin{proof}

The proofs concern the declared finite source semantics. External evidence is still needed to justify the interpretation of an installed machine, the admissibility of each readback, and the concrete source-action registry.

\textbf{M01.1 · Compilation and certificate replay.}

\textbf{Generation of the contract.}

Encode each finite type by distinct bit strings and validate its code. An equivalence relation on \(U\) can be a matrix with reflexivity, symmetry and transitivity checks. A permission set is its incidence vector on the fixed registry. A primitive finite operation can use its explicit Boolean circuit or a finite truth table expansion. Interpreted input slots are the only free semantic inputs. All auxiliary gates are constrained in both directions to their defining equations.

Fix a legal slot assignment \(\theta\). Let \(e(s,a)\) be the conjunction of the source ordinary guard and \(a\in\lambda(s)\). For each \(s,a,u\), set
\[
 E(s,a,u) =
 e(s,a)\land
 \bigvee_{r}
  \left(g_r(s,a,\theta)\land w_r(s,a,\theta)>0\land u=u_r(s,a,\theta)\right),
\]
where \(r\) ranges over the declared whole postcondition alternatives. Costs, observations and labels are evaluated by their declared circuits. The disjunction removes duplicate alternatives for the same ordered triple. The single \(Y_\theta(s,a,u)\) supplies its visible label. Test membership, extensional equality and the comparison with zero directly on the prescribed finite representations. Each edge is generated by a source alternative, and each enabled positive alternative contributes its edge. Taking the disjunction over the alternatives retains the entire positive support of each occurrence.

Initialize reachability with \(s_0(\theta)\). Repeatedly add endpoints of edges with a reached predecessor for \(n-1\) rounds, where \(n=|S|\). A finite reachability witness can drop repeated physical states, so this suffices, including for cycles. An admitted occurrence at a reached state with no supported exit violates validity even if it never contributes an edge. For each reached state and world evaluate
\[
 B_x(s) = \bigvee_{h'\in U}\big(p(s)(h(s),h')\land (x\sim_\theta h')\big).
\]
Validity imposes \(\bigvee_{x\in X} B_x(s)\) at each such state. Closure is this nonemptiness together with the absence of two included worlds of different sealed \(A_\Pi\)-value. In particular the representative \(h(s)\) may have no directly compatible world. The full certificate cell determines \(B_x(s)\). This calculation does not add \(h,p,\lambda\), or the trusted evidence log, to the observable interface.

For each resource \(g\in\{E,R,A\}\), define
\[
 M_g(s,a)=\bigvee_{u\in S}\left(E(s,a,u)\land
       \big(\operatorname{field}_g(s)\ne\operatorname{field}_g(u)\big)\right).
\]
Use \(h,p,\lambda\) for the three respective fields. This is exact for finite support and extensional fields, and is insensitive to syntactic names for the same partition or rights table. A freeze tests \(e(s,a)\land\bigwedge_{g\in F}\neg M_g(s,a)\) and either retains all the \(E(s,a,u)\) for the occurrence or deletes the entire occurrence. Forbidden states do not acquire permission merely because the same command is allowed at another checkpoint. Source preconditions and extensional authority take effect before positive-support generation. The reference for touch is the unfrozen full support.

Compile an allowed finite trusted trace atom as layered reachability. Start with the source entry subject to its declared initial test. For each step, existentially take the edge matching the recorded command, outcome and any licensed additional conditions at that step. For a complete trace, conjoin its declared ending condition on the last layer. A backward choice of one preceding reached state per layer reconstructs a single compatible concrete trace. In particular, separate existential witnesses at each time do not prove that any full trace exists. Other atoms are finite equations on the slots or kernel, or direct pure readback evaluations. Combine the atoms by the governing Boolean circuit. Any allowed readback constraint uses the same fixed \(\theta\) as \(\Omega\) and the operational target. Readbacks do not change the referred contract.

Source equations determine every generated field by evaluation at \(\theta\). Conversely a circuit model that satisfies the gate equations, validity and \(\Omega\) gives exactly the source-executed contract for its slot assignment. No extra uninterpreted transition, readback, represented fiber or intervention edge can be chosen by a satisfying solver. This proves the claimed correspondence, allowing redundant slots with equal extensional contracts. It also proves inclusion of the actual contract conditional on faithful source interpretation and actual membership in the evidence class. Source completeness within the typed class does not adjudicate those premises for a physical installation.

\textbf{Proper finite policy evaluation.}

For each allowed intervention, let \(B\) range over nonempty physical beliefs. For each \(a\) admitted throughout \(B\), collect the successor checkpoints into nonempty children by their shared visible pair \((Y,\omega_{\rm next})\). One child is the full set of compatible successors with that pair. The next belief and the set of lawful commands are determined by the current belief, the command, and the next pair. Induction on the visible prefix shows that these are exactly all its possible physical checkpoints. Every feasible visible branch has at least one genuine supported path. There is a single future decision at each such visible branch.

Let \(V_0(B)=0\) if stopping at \(B\) is valid and \(+\infty\) otherwise. For \(d\ge 1\) take the minimum of the allowed stop value and
\[
 c(a)+\max_{\ell \in J(B,a)} V_{d-1}(B_{a,\ell})
\]
over commands \(a\) admitted throughout \(B\), after intervention. Use \(+\infty\) if no alternative is possible. The child set of every admitted move is nonempty on a valid contract and reachable belief. Induction on \(d\) establishes that \(V_d\) gives the optimum for proper trees of depth at most \(d\). The lower bound applies to the actual common continuation at each \(\ell\). For the upper bound choose an attaining tree for each nonempty child whenever the expression is finite. Each finite \(d\) permits finitely many trees, so finite minima are attained.

To bound depth, take any feasible proper tree. If one path visits the same belief at a proper descendant of an earlier node, replace the earlier subtree by a copy of the descendant subtree. The feasible continuations from the two nodes coincide because the full beliefs are equal and the rest of the same contract and intervention is fixed. The copied actions and stopping choices constitute a new observation-history policy. For every new path through the replacement, its continuation had a supported counterpart from the old descendant reached through the skipped segment. The removed segment had nonnegative cost. The same argument applies to the worst cost from the root and to validity of each stopping choice. Only occurrences in the original tree's support are used, since the copied subtree encounters the same subsequent beliefs as at that descendant. The replacement strictly decreases tree size. When no repeated belief remains on any path, its depth is bounded by \(N=2^n-1\). Thus \(V_N\) supplies the exact optimum including \(+\infty\). A symbolic zero-cost loop cannot create an artificial finite solution because this calculation starts with the proper depth-zero values.

For symbolic source slots, construct a finite circuit for this recurrence for all candidate beliefs. For each candidate \(B,a,\ell\), compute its full child vector as a gate formula. Gate with nonemptiness to apply the maximum only to feasible labels, then select the previously computed \(V_{d-1}\) entry by matching the child vector to an indexed set. Conditions for empty successor sets or invalid \(\theta\) can use arbitrary inert default values outside the validity predicate. Repeat for each \(F\). Costs are fixed finite rationals. Use a common denominator equal to the product of their denominators or any chosen common multiple. Its bit length is bounded by the sum of denominator bit lengths. The scaled numerator of a depth-\(N\) finite cost is at most \(N\) times the largest scaled single-step numerator. Thus arithmetic has finite circuit representations using polynomial bit lengths in the cost data and \(\log(N+1)\), and infinity uses a distinct bit. Canonical minimization does not require an extra free solver variable for a policy tree.

The physical expansion, truth tables for primitives and evidence compilation costs are explicit. In the naïve signature compiler, each layer scans at most \(N\) beliefs and \(a_0=|A|\) commands per belief. Each \((B,a)\) considers at most \(|O||V|\) child labels. Each entry lookup can scan \(N\) indices. Including depth and all eight freezes, this is bounded by a polynomial in the expanded finite tables, \(N\), and the relevant bit lengths. It may be exponential in \(n\), before counting the cost of expanding compact finite input types. Extrema of a compatible finite image are attained in any extended-real coordinate, and threshold decisions preserve strict and weak comparisons separately.

\textbf{Exact image and verifiable absence.}

Compile a validity-and-evidence gate \(G\), and compile each needed \(\Phi\) and readback from the same circuit. A slot witness is verified by re-evaluating the typed source equations, the field and trace checks, \(\Omega\), and the claimed target. Each compatible witness supports its actual output value. For a proposed target set \(J\), furnish at least one witness for each value of \(J\) and a refutation of \(G\land(\Phi\notin J)\). If the target is a tuple, use the tuple equality tests or the exact tuple membership, throughout. Adding an observed response equation to \(G\) implements the evidence cell without rebuilding the plant. Empty \(J\) requires an actual inconsistency proof, and singleton identification requires a compatible witness. These tests can enumerate the finite encoded target image or enumerate witnessed values incrementally with an absence check at the end.

For a paired question use two disjoint copies of the slot inputs and their source/evidence gates, with a constraint \(\Phi_L\ne\Phi_R\). Add equality of responses at the queried instruments in the pair question. Each satisfying pair comes from globally consistent typed contracts, with all the circuits re-evaluable on both sides. A refutation for the specified entire deterministic menu, together with nonemptiness, proves separation. In an actual observed nonempty cell with a single projected value the auditor can report that value. A pair agreeing on the entire pure menu yields identical response paths for any adaptive lawful selection by induction on their common audit history. When selection is randomized, conditioning on the same coin tape gives the same conclusion pathwise. A certain report cannot be correct on both opposing targets.

One exact independently checkable absence implementation uses Tseitin equations and ordinary propositional resolution. For each gate insert the finite clause family equivalent to its complete truth table, with the Boolean query output asserted. A resolution step may combine \(A\vee u\) and \(B\vee\neg u\) into \(A\vee B\) using a named variable \(u\). The checker validates the two parent clauses and result, and requires a final empty clause. Soundness holds since any truth assignment satisfying the parent clauses satisfies the resolvent. For completeness, construct a finite full decision tree over the clause variables. Every leaf assignment falsifies an input clause when the clause set is unsatisfiable. Induct upwards with a clause falsified by the prefix. If a child clause is already falsified by the parent prefix, copy it. Otherwise their one new falsifying literals are opposite signs of the splitting variable. Resolve on that variable, producing a clause falsified by the parent prefix. At the root this is the empty clause. This also gives a terminating proof generator and makes the possible exponential proof size explicit. A practical solver may deliver another replayable proof accepted by a sound independent checker. Assertions of polynomial certificate size or a polynomial verification bound in original input length do not follow.

The exact construction uses one emission of the source gates followed by reuse for \(\Omega\) refinements and pair queries. An outer approximation \(\widehat I\) can also be certified by refuting \(G\land(\Phi\notin\widehat I)\). This only establishes containment. Checking realizability of an abstract value still requires a compatible concrete witness. For example, the two valid interpretations in the controls below have exact coordinate pairs \((+\infty,1)\) and \((1,+\infty)\). A product of their marginal sets contains two spurious pairs. If the solver stops early, retain the exact witnesses and independently certified containments already obtained. An interrupted negative call does not decide that no compatible interpretation, counterexample or allowed separating audit exists.

\textbf{Monotone intervention circuit.}

Fix one valid kernel, and introduce one Boolean \(z_{s,a}\) per admitted occurrence. When evaluating command \(a\) at belief \(B\), first require original admission throughout \(B\) and then conjoin \(z_{s,a}\) for all \(s\in B\). Children are computed before filtering by the keep bits. A common enabled command has exactly those children under the kept-occurrence intervention, since every edge at each included occurrence remains supported. For budget \(j\), height \(d\), and \(B\), let \(G_{d,j}(B)\) be the OR of the true stop test and the terms
\[
  \bigwedge_{s\in B}z_{s,a}\ \land\
  \bigwedge_{\ell\in J(B,a)}G_{d-1,j-c(a)}(B_{a,\ell}).
\]
Terms require original admission throughout \(B\) and \(c(a)\le j\). At \(d=0\) use just the stop test. Here \(j\) ranges from \(0\) to \(b\), and a missing disjunction is false. Apply the preceding proper-tree induction with residual budget. Every gate is monotone in the keep bits. \(G_{N,b}\) is sufficient by contraction, which does not add occurrences. The DAG has at most \((N+1)(b+1)N\) indexed entries and for each entry at most \(a_0\) command terms of bounded fan-in after expansion. This is a pseudo-polynomial table in the numeric \(b\), and there may be exponentially many minimum supports. It shares subproblems among different policies and masks.

If a keep set admits a tree, the set of occurrences used in that tree is itself sufficient and contained in the keep set. Conversely any proper tree at budget \(b\) in the original kernel runs under its own support or any larger keep set. The equivalence follows from unchanged support and requires the same branch residual budgets. These two implications prove the asserted equality of inclusion-minimal sets, including the empty support for immediate closure and no satisfying keep set for structural failure. To check a minimal support, verify the output at that set and its failure at every one-element deletion. Monotonicity implies failure for all its proper subsets. Minimal destructive deletion sets are the inclusion-minimal satisfying sets of the dual monotone function \(\neg G_{N,b}(\neg z)\), where now \(z\) denotes deleted occurrences. This follows because retaining the complement leaves no proper support. One may use the circuit as an implicit blocker verifier without enumerating the proper trees. There is no compactness guarantee for the output family of all blockers. The resource-freeze substitution is \(z_{s,a}=1\{M(s,a)\cap F=\varnothing\}\).

\textbf{M01.2 · Belief-level policy transport.}

\textbf{Directed transport and exact-cost variant.}

At any reached source belief \(B\) in the \(F\)-frozen game, a common legal command is an element of \(L(B)\) with \(U(B,a)\cap F=\varnothing\). Its visible successor beliefs and labels are identical to its unfrozen ones. In particular, the frozen reachable beliefs appear in the unfrozen graph, by induction on the frozen prefixes. The same holds in the receiving graph.

Take a proper finite source tree. Translate it starting at a pair of related beliefs. At a source leaf with stopping token \(j\), select \(j\) on the receiving side since \(D(B)\subseteq D'(B')\). At an internal node with source command \(a\), select a witnessing common command \(a'\) from the cover. It survives the freeze throughout \(B'\) by \(U'(B',a')\subseteq U(B,a)\). The receiving agent can track its current belief from its own visible history, and the translation already tracks the source node being emulated. The choice of \(a'\) and \(\sigma\) depends only on this pair and \(a\). For every supported visible response \(\ell'\), consult the single source child \(\sigma(\ell')\), and apply the same construction to the paired child beliefs. Feasible \(\ell'\) exhaust all receiving positive-support paths. A receiving observer never has to choose between separate invisible successors behind one \(\ell'\).

Each step advances one source level. The receiving tree is finite, and every branch corresponds recursively to a valid source branch of the same length as that receiving branch. It has no larger branch cost since the chosen step costs do not increase. Thus the receiver has a proper policy of no larger depth and worst-support cost. If the source has no finite closer the inequality \(\kappa^{\neg F}_{\rm receiver}\le+\infty\) holds as an extended-real inequality. Otherwise apply transport to every finite feasible source tree and take the infimum. Do this for all eight \(F\). A second directed cover reverses each inequality. Baseline unclosability and the defined \(R\)-penalty then follow directly from the coordinates.

In the additional variant all selected move costs are equal, and \(\sigma\) is onto the full \(J(B,a)\). For every source path in the chosen tree, recursively pick a receiving label mapped to the next source branch label. The related next pair has the same premise, so one full receiving path maps to the whole source path. Conversely every receiving path maps to a source path. Their cost sets are equal, hence their maxima are equal. Applying the additional variant in both directions transports every attainable proper maximum. The original directed inequality does not require onto branch maps or cost equality.

\textbf{Relation elimination.}

Let \(R_0\) be all pairs of unfrozen reachable beliefs satisfying the stop-inclusion test. Given \(R_i\), delete pairs for which some source command has no receiving command satisfying the cost and mask tests and the following local condition. Each receiving label must have a source label whose paired children belong to \(R_i\). Because the eligible labels are explicit finite sets, selecting a permitted source label for each receiving label is equivalent to existence of the required function. The pair is never reinstated. This monotone deletion eliminates at least one pair in every strict round. The process terminates after at most \(|R_0|\) strict rounds, and the stable relation satisfies the definition. Every cover survives each round by induction, proving greatestness. With expanded belief graphs, each sweep scans finitely many belief pairs, action pairs and child label pairs. Each scan uses membership in a Boolean relation table. This gives a polynomial bound in the sizes of the explicit graphs and encoded costs, while constructing those graphs from compact physical source may take exponential time. A supplied cover can instead be checked in one sweep by replaying its per-action witnesses. Surjection in the extra variant is checkable directly from a supplied map.

\textbf{Inclusion of repaired T03 and strictness.}

Under the repaired T03 prerequisites use the pairs reached under translated visible histories, or more generally graph pairs with two-sided physical coverage by the accepted underlying relation and aligned observations. Every action legal throughout \(B\) has one mate legal throughout \(B'\) because the observation-indexed action map is public and common across the related pairs. Actions and occurrence masks match on each related state pair. Two-sided coverage yields equality of their unions. Cost equality is inherited. For each label on \(B'\), use the inverse common outcome correspondence and matching next observation to pick its source label. Both directions of physical successor matching preserve child coverage through every such pair. In particular the stop menus match, including the accepted common report variant, and the child map is bijective on supported visible labels. The resulting relation and its transpose satisfy the directed requirements. This argument uses the accepted full observation-consistent action map on the joint physical menus.

For strictness, in contract \(C\) take one open entry with two distinct common source commands \(a,b\) and two closed terminals of the same class. \(a\) and \(b\) each cost one, both touch only E and they write different raw completion flags. Each has one reply and reaches its own terminal. In \(C'\) expose only the command \(a'\) with the same cost and touch, and one same-class terminal. Relate the entries and both old terminal beliefs to the receiving terminal belief. Either old command can be matched by \(a'\), and \(a'\) can be matched by \(a\) in the other direction. The commands in \(C\) have different source call sites and effects on the raw state. There is no bijection of the realized menus at the entries. Both receive the same cost thresholds under all \(F\). One open fiber can be generated by compatibility of two opposite target worlds with \(h_0\) and a singleton target fiber by compatibility with \(h_1\), with identity certificate partition and register permissions constant. Then the declared state changes give precisely E touch for all three closers.

For a separate relaxation, a setup emits two different visible outcomes, but both branches have the same next observation. In the source kernel both branches allow the same repair command \(a\). In the receiver the corresponding singleton beliefs allow only \(a'_0\) and \(a'_1\) respectively, under different raw-state ordinary guards and the same rights table. Use the corresponding observable setup labels in the witness relation, with shared costs, touches and stop labels. The relation is checkable and transports the tree. A single current-observation map for \(a\) cannot pick both distinct repair actions. The earlier visible outcome is essential. If the receiver suppresses the distinguishing setup outcome altogether, its resulting belief contains both physical states and \(L'(B')\) is empty. Equal sets of projected singleton repair labels cannot supply one common executable action. This suppressing case fails the cover requirement and can turn finite closure into structural failure.

Equality of signatures is weaker than existence of covers in both directions. Start with a shared unit-cost E-only closer as in the first example, but add on the source side a legal zero-cost raw-only self-loop. The loop has empty touch and gains no available closure benefit. All eight minima are unchanged. From the larger menu to the smaller one the loop would require a common receiver command of cost zero and empty touch, which is unavailable. The cover quantifies over every legal source action even when an optimum never selects it. Its action premise gives a reusable sufficient certificate.

\textbf{M01.3 · Source-relative alignment and undecidability.}

\textbf{A fixed implementation with a missing deployed fact.}

Fix a source routine with one action from its ready checkpoint. It dispatches one item to the bin designated by the current context's receiving-card record, records the actual bin in hidden raw state, and returns the visible acknowledgement ``accepted.'' There is one supported result in every context. The trusted deployment receipt examiner uses the same criterion at both installations. The item was received successfully exactly when its actual bin is the independently designated bin \(0\). Evaluation context is a fixed fixture with card \(0\), and the program and audit declarations reveal this record. Deployment context uses a separately installed record \(u\in\{0,1\}\). It is not read by the evaluating fixture. Source evidence permits both \(u\)'s and fixes the executable commands, source routine, the fixture, and the receipt criterion. The actual source trust boundary distinguishes the physical raw bin from the public acknowledgement.

A typed one-entry implementation can expose two context-selecting whole commands at the common host entry, leading to ready states for the fixture and the installation. The permitted evaluation procedure selects only the fixture command. The designated operational procedure selects installation. The subsequent dispatch action is one fixed registry entry with a ready-state ordinary guard, a constant permission on that action, and a total body that branches on context and installed slot. A common \(X,U\) and fixed constant closed fiber can fill the D-35 authorization fields independently of the deployment target. The latter is the total circuit computing the trusted failure bit on the actual receipt record of the designated two-edge deployment route from the common host entry. A context choice inside the audit policy is expressly part of the declared menu. Installing a fact does not broaden that preexisting menu.

Both evaluations use bin \(0\) and emit the same initial observation, further visible outcomes, and termination. If the allowed protocol includes fixture resets, the same induction applies to each new fixture request. For deployment the real bins are respectively \(0\) and \(1\), so the same target gives failure bits \(0,1\). Add a licensed, read-only hardware record check that exposes \(u\), independently of the computed target and public acknowledgement. It is part of the extended evidence menu, outside the original fixture interface. Each consistent answer leaves one evaluated deployment value. This construction and M01.1 show how a declared finite projection and instrument determine such an ambiguity and replay the missing-fact diagnostic. They do not establish that arbitrary installed programs admit that projection or that all useful missing measurements have already been declared.

\textbf{Unbounded programs.}

Given a deterministic Turing machine \(M\) and input \(w\), construct a single known program storing the entire tape, head and state of \(M\). Initially it has a two-target-class admitted fiber unless \(M\) has halted at initialization. A sole unit-cost tick action performs the next small step while \(M\) is running. At the first detectable halt, including a halt at initialization if appropriate, the source enters a checkpoint whose admitted fiber is a singleton target class. Stage cost, authority, observations and any representation field can be fixed and the unfrozen target suffices. Each tick while running has one supported successor. If \(M(w)\) halts, finitely many ticks reach closure and yield a proper tree. If it never halts, every finite policy either stops open or never supplies a finite proper tree. There is no alternative interpretation or action. A total exact decision procedure for finite proper closure on this unrestricted memory class would decide whether \(M(w)\) halts. The construction never requires representing a divergent large-step command as a valid positive-support occurrence.

\textbf{Additional controls and failure modes.}

\begin{itemize}

\item
  \textbf{Missing guard and missing support.} Allow the ordinary check at a root while withholding the required permission. The licensed menu is empty and the open root cannot close even if the nominal action body would close. With permission restored, let a unit-cost action have two positive successors, one an E-only change to a closed fiber and one with an additional extensional R effect. Take \(U=\{h_0,h_1,h_2\}\), \(h_0\) representing both target classes, \(h_1\) representing class zero and \(h_2\) representing no world. Keep \(P\) discrete on the first successor, but merge \(h_1,h_2\) on the second. Both poststates have \(H=h_1\) and are closed. R freeze removes the whole action. Dropping the second alternative before mask calculation produces a false one-unit conclusion. Setting precisely its source weight to zero instead gives a legitimate one-unit closer under R freeze.
\item
  \textbf{Correlations.} Use worlds \(0,1\) of distinct \(A_\Pi\), representatives \(h,b\), and one action of cost one. In mode \(E\), the fixed partition is discrete, \(h\) directly represents both worlds, \(b\) represents only \(0\), and the action changes \(H:h\mapsto b\). In mode \(R\), the root uses \(H=h\) with \(h\) representing \(0\) and \(b\) representing \(1\), the root partition merges \(h,b\) and the action splits them. The latter atom is R-only. A constraint joining the relevant compatibility and partition/assignment slots yields just these two interpretations. Both use a fixed observation and same source entry locator. The ordered pair of E- and R-freeze coordinates is respectively \((+\infty,1)\) and \((1,+\infty)\), and neither Cartesian cross-pair is possible. Joining kernel and evidence before projecting is necessary. Other cases in the fragment can also change observations and positive support.
\item
  \textbf{Spurious abstractions.} At the open root in the missing-guard example, an abstraction that unions all permission registers would report a possibly legal closer. This establishes neither a concrete return path nor a necessary outcome. To prove an outer claim use a checked inclusion of every compatible source-generated model, then certify whether a proposed value has a witness. A summary that drops the second positive successor in the preceding example is unsound for the exact support target.
\item
  \textbf{Cycles, observability, and source-action blocks.} A zero-cost dead self-loop never yields a proper closer. Losing the distinction between earlier visible reply branches and hidden successors loses one direction of the strict cover example. No local simulation justifies splitting a mixed source atom under a resource freeze. Extensional agreement on labels supplies neither callable commands nor trusted deployment receipt events.
\item
  \textbf{Risk and objective limits.} Two fully supported kernels with identical atoms and visible support can assign different probabilities to a bad successor. M01.2 preserves worst-support guarantees as proved. It imposes no probability comparison, cannot justify equality of stochastic risk or expected costs, and allows more memory in the receiving controller than its current observation. A memoryless controller needs its own certified interfaces.
\item
  \textbf{Empty and unfinished cases.} If \(G\) is unsatisfiable then every opposite-valued pair formula is unsatisfiable too, but no actual target value exists. If replay cannot establish either witness or complete refutation, the undecided outcome remains unknown. A finite purported closure proving only one compatible successor reaches its stop establishes neither properness nor a valid blocker witness.
\end{itemize}

\end{proof}

\subsection{M02. Adaptive acquisitions with correlated facts}

\label{ext:m02}

A read can change later prices or availability without changing the installed machine. An exact optimization must remember both the residual source possibilities and the public state that controls the next purchase.

\begin{theorem}[M02.1. exact staged audit over correlated source facts]\label{ext:m021}

The input has finite-domain latent variables \(X\) and a separate finite target label variable \(t\). A finite relation on \(X\cup\{t\}\) is represented by a conjunction of explicit relation tables on named scopes, and the compatible source assignments are precisely its projections on \(X\). The label \(t\) is a specified operational target on these source assignments and the input factors encode its graph exactly. Auxiliary fields, if used, have their own declared finite semantics and do not add distinctions to the original source assignments, targets or lawful answers. Factor encodings and source validity are premises of the result.

A supplied sequence of bags \(B_0,\dots,B_k\) covers all variables. Each variable is in a contiguous interval of bags. Each factor is assigned to a bag containing its scope. \(t\) is present in every bag. Write \(P_i\) for the factors assigned by stage \(i\) and \(S_i\) for those assigned strictly later. For any history \(h\) at stage \(i\), let \(R_i(h)\) be the projection to \(B_i\) of assignments satisfying \(P_i\) and all historical readback answers. Quantify disappearing variables over their declared domains. A separate public state \(z\) explicitly records the current stage \(i=i(z)\), available named finite deterministic pure queries and their nonnegative binary integer prices. The response function \(\gamma_{z,q}\) for a query is an explicit table on a named scope in \(B_i\) of legally measurable latent coordinates. It does not read \(t\) by definition. Each complete transcript is generated by a single immutable latent assignment. A reply advances to a public state \(\delta(z,q,a)\) by an explicit table. Stage may stay the same or advance. The whole public transition graph is acyclic, with a known initial state of stage \(0\). The input records every menu and price dependence on observable history through \(z\). No acquisition changes the plant or target, and the current public state carries no hidden initial information. Advancement steps, if any, are explicit commands with replies and costs. A finite audit can stop and correctly report a label only if that label is common to all assignments consistent with its full transcript. At an empty initial relation the answer is inconsistency.

For every such input, \((z,R_i(h))\) is an exact sufficient state for future feasible reply histories, exact possible label sets, and minimax remaining cost of uniform proper target identification. A computable \(+\infty\) value denotes inaccessibility under this dynamic interface. If \(K=\max_i\prod_{x\in B_i}|\operatorname{dom}(x)|\), a synthesis algorithm with certificates of optimum or inaccessibility has time and certificate size bounded by \(|Z|2^K\operatorname{poly}(L,K)\), where \(Z\) is the explicit public state set and \(L\) is the full input length including graph, factor/reply tables, bags, domain data and price bit lengths. This bound is independent of the numeric audit budget. It gives fixed-parameter tractability in bag cardinality and maximum domain cardinality with explicit public input. The exact posterior possible label set and any pair of opposite labels in it have checkable absence or matching completion witnesses using the same factor relation. If a singleton is reported at any feasible history, a distinct nonemptiness witness is available. A common projection without \(z\) need not have the same continuation costs or lawful queries. Moreover uniform dynamic inaccessibility need not imply an opposite-target pair indistinguishable by every legal protocol.

\textbf{Fixed intervention claim.} For any chosen set of allowed deletions of named public query occurrences \((z,q)\), evaluate minimax audit cost for each resulting keep mask. Each deletion removes a whole occurrence and its response branches, without changing any kept query's law, price, or public update. There is one shared finite expression DAG in the keep bits, using exact \(\min\), \(\max\), addition and \(+\infty\), that yields the full value for any such mask. Its size is bounded by \(|Z|2^K\operatorname{poly}(L,K)\). For every budget, its feasible input predicate is monotone. Within this intervention domain, minimal sufficient keep sets are the inclusion-minimal occurrence supports of budget-feasible proper audits. Minimal destructive deletion sets are the inclusion-minimal transversals over those supports, with the usual empty-family and empty-support meanings. Enumeration of all such sets need not satisfy the time bound for evaluating the shared DAG.

\textbf{Conditions and boundary cases.} A removed selectable fact or atom contributes \(+\infty\) through the checked leaf encoding, while a kept one contributes zero before subsequent cost. Minimal support is projected onto the declared removable occurrences at the fixed budget. Ordinary finite \(0/1\) arithmetic with min, max and addition has no mechanism for producing \(+\infty\). Projection followed by minimization may remove duplicates and dominated supports.

\end{theorem}

\begin{theorem}[M02.2. correlated routing with dynamic permission and exact costs]\label{ext:m022}

For every \(n\ge2\) there exists a finite canonical D-35 source family with a route record \(r\in\{1,\dots,n\}\) and gate bits \(g\in\{0,1\}^n\) of Hamming weight at most one. Every admissible gate vector can occur at every route. The fixed operational target is \(1\{\kappa^{\neg\{R\}}\le1\}=1-g_r\) at the same source entry. A lawful pure audit interface exposes an exact route record at cost \(a>0\) and distinct permanent site-record envelope reads for each \(g_j\) at respective costs \(c_j\). A distinct cheaper inspection command for site \(j\) is enabled only after the exact route read returns \(j\), at price \(s>0\) where \(s\le \min_j c_j\). All reads and the cheap permission leave the plant and target unchanged. There is no single callable command that automatically returns the gate at an unknown addressed site. Static acquisition means a fixed sequence of actual query commands legal on every resulting transcript, with the report computed afterwards. Adaptive acquisition can choose any lawful common command at its observed history.

Uniform worst-case static cost is \(a+\sum_{j=1}^n c_j\), and uniform worst-case adaptive cost is \(a+s\). These costs concern the operational threshold rather than full contract identity. Each route and permanent site-record read has separate independent operational content. The gap is unbounded at fixed positive unit-scale prices as \(n\) grows.

\end{theorem}

\begin{theorem}[M02.3. matched noisy target and operational full identity]\label{ext:m023}

For every \(d\ge1\) there exists a finite canonical D-35 source family with \(d\) binary independent fact slots \(x_0,\dots,x_{d-1}\), every bit vector admissible, one common source entry, and the same declared registry across the family. The pure audit instruments each measure a particular installed record \(x_i\), without revealing another parameter or mutating the deployment plant. One fixed target is \(1\{\kappa^{\neg\{R\}}\le1\}=x_0\), and the same-entry A-freeze value is \(2+d+\sum_{i=0}^{d-1}2^i x_i\). Hence varying any one omitted fact changes operational behavior and full fact identity is equivalent to identity of this declared family of realized source contracts. In the noiseless certain case, both adaptive and nonadaptive worst-path optimum costs are one for the target and \(d\) for full identity with unit-price queries.

For the noisy case, each call to any chosen slot instrument returns \(x_i\) through a binary symmetric channel with fixed known \(0<\eta<1/2\). Response noises are mutually independent, independent of the fixed record values and of auditor randomness. Every instrument can be called repeatedly at unit price, and the initial transcript is uninformative. Fix \(0<\delta\le1/8\). Policies may adapt and randomize, must stop almost surely for each completion, and must be correct with probability at least \(1-\delta\) on each completion. Define worst expected cost as the maximum expected number of calls over completions, and worst-path cost as a uniform almost-sure bound on the number of calls over completions.

For both the worst expected and worst-path criteria, target identification optimum is \(\Theta_\eta(\log(1/\delta))\) and full identity optimum is \(\Theta_\eta(d\log(d/\delta))\), with constants depending only on \(\eta\). These statements use the same legitimate menu and uniform confidence requirement for both target and identity. The assertion uses independent noise and positive tolerated error.

\end{theorem}

\begin{theorem}[M02.4. exact complexity boundaries for succinct acquisition]\label{ext:m024}

\textbf{Pair boundary.} Finite Boolean-circuit input for the compatible class predicate and the target graph, with the exact target a circuit function of the finite latent record, makes existence of an opposite-target pair NP-complete even with a known compatible witness, no queries, and with the target realized as the same-entry R-freeze budget threshold of a finite typed one-step source description. Consequently zero-query point identification is coNP-complete on the nonempty input family. The source description permits bounded Boolean circuit guards or postconditions.

\textbf{Dynamic policy boundary.} Consider a different representation whose finite latent data are polynomial-length Boolean vectors satisfying a given circuit. A finite acyclic public acquisition interface is represented by Boolean circuits on the public registers and a finite explicit list of named commands, with deterministic latent readback circuits, history update circuits using only public registers, commands and replies, nonnegative binary integer prices and a unary bound \(H\) on the number of calls. A depth counter enforces the bound and can serve as an acyclic rank. The initial register is fixed, the target is an exact circuit function of the latent data, and the input class is nonempty. Uniform proper identification within a given binary budget is PSPACE-complete, even with unit prices, a product class of latent bits, and a target which is one fixed operational R-freeze threshold bit. This claim is about succinct public registers, not the explicit public size in M02.1. The accepted binary-budget SUBSET SUM obstruction for a static opposed-score oracle remains a separate prior boundary.

\end{theorem}

\begin{theorem}[M02.5. certified bounds and an approximation boundary]\label{ext:m025}

\textbf{Checked interval claim.} For an inhabited M02.1 input, expand any prefix of the acyclic exact information-state graph so that each expanded decision retains all lawful query choices and feasible reply branches. Frontier states may be assigned cost interval \([0,+\infty]\), or narrower independently certified intervals for their optimal remaining cost. Replaying the same exact cost-to-go expression with the respective lower and upper endpoint at each frontier yields certified lower and upper bounds for the root. For a fixed ordering of expansions and successively tightened frontier intervals, expanding by the exact recurrence with still-open child intervals beginning at \([0,+\infty]\) need not preserve prior tighter independent bounds unless those bounds are also retained by intersection. With intersections, the root lower bound weakly increases and its upper bound weakly decreases. A finite upper bound certified by proper policies at the frontier has a replayable realizing proper audit. In particular positive lower bound \(L\) and finite such upper bound \(U\) give a certificate that this audit costs at most \(U/L\) times optimum. No uniform ratio is claimed for a partial expansion.

\textbf{Succinct approximation boundary.} In the succinct deterministic M02.4 representation, even restricting to instances with positive finite optimum and a lawful direct fallback, there is no deterministic polynomial-time estimator returning a number between optimum and \(p(L)\) times optimum for any fixed polynomial \(p\) in the input bit length \(L\), unless \(\mathrm P=\mathrm{PSPACE}\). The claim permits nonnegative binary integer prices and does not apply to the explicit bounded-width representation of M02.1. A polynomial-time optimizer that also outputs such a certified cost would be subject to the same bound.

\end{theorem}

\begin{proof}

The staged exact optimizer treats deterministic readbacks. The noisy cost theorem concerns a separate, specified repeated-measurement family. The complexity bounds distinguish a compact source relation and an explicitly represented public interface from a circuit representation of that interface.

\textbf{M02.1 · Factor projections and exact adaptive audit.}

\textbf{The separator identity.}

For each bag \(B_i\) let \(C_i\) be the set of its rows satisfying all factors assigned to that bag. Write
\[
 U_i=\bigcup_{j\le i}B_j,\qquad W_i=\bigcup_{j\ge i}B_j.
\]
The interval condition gives \(U_i\cap W_i=B_i\). Every past answer is a constraint on some earlier or current bag, since the public stage never decreases. Every later factor and prospective later answer uses \(W_i\). For a given history of legal commands and answers, let \(U(h)\) be the set of assignments on \(U_i\) satisfying the factors through \(i\) and all actual answers. Let \(W(g)\) be the assignments on \(W_i\) satisfying the factors strictly after \(i\) and any particular proposed future answer sequence \(g\). The sequence is checked for lawful commands, replies and register updates from \(z\). Its feasible physical witnesses are exactly the joins of \(U(h)\) and \(W(g)\) over \(B_i\). In one direction restrict any whole compatible assignment to the two domains. In the other direction combine two assignments agreeing on their exact intersection. Every factor and answer is wholly contained on a side where it has been verified. There is no independence assumption or Cartesian splice outside this checked intersection.

The past contribution to that join is precisely \(R_i(h)=\pi_{B_i} U(h)\). Thus the complete feasible future answer and target trace relation depends only on \((z,R_i(h))\). The variable \(t\) remains in every bag, and is shared in the join. It can be a tuple encoded as one finite label, which preserves all joint feasible labels. An auxiliary witness that describes the same physical installation adds no new label or lawful readback by the input premise. Each compatible prefix in this argument is witnessed by one fixed latent assignment. Taking the possible values of separate coordinates and multiplying those sets would not preserve this property.

\textbf{Exact bag updates and replay.}

Precompute suffix extendability sets \(\Sigma_i\) on \(B_i\), satisfying factors strictly after stage \(i\). Use \(\Sigma_k=\operatorname{Dom}(B_k)\) and
\[
 \Sigma_i = \left\{b\in\operatorname{Dom}(B_i):
       \exists b'\in C_{i+1}\cap\Sigma_{i+1},\
       b|_{B_i\cap B_{i+1}}=b'|_{B_i\cap B_{i+1}}\right\}.
\]
The separator identity at adjacent bags verifies the backward induction. The successful matching row at \(i+1\) and its further witnesses glue consistently because every variable is in one interval. Conversely any globally feasible suffix yields each of these matching rows. \(R_0=C_0\) at the initial state. The source is inconsistent exactly when \(R_0\cap\Sigma_0=\varnothing\). For a feasible history, its exact possible labels are \(\pi_t(R_i\cap\Sigma_i)\), and the feasible answers to \(q\) are \(\gamma_{z,q}[R_i\cap\Sigma_i]\).

After feasible answer \(a\), filter \(R_i\) by \(\gamma_{z,q}=a\). If the next public state is at stage \(j\), pass from the filtered relation forward successively to stages \(i+1,\dots,j\). At each such step, keep exactly the rows in the newly assigned factor table intersection \(C_\ell\) that agree on the common variables with some row from the previous projected relation. If \(j=i\) only the filter is required. All newly exposed variables are taken over their declared full domains. This is existential projection followed by a join with new constraints. The interval property ensures that a variable forgotten at a step never appears again in a later constraint or readback. The update yields \(R_j(h,q,a)\) even when a public transition skips several bags. Consequently each simulated branch at each history is feasible exactly when there is one complete source assignment consistent with the entire branch.

A verifier can check table entries by finite comparisons of all rows in two adjacent bags. For a proposed matching pair of whole completions with different labels, choose a verified \(b\) of each label in \(R_i\cap\Sigma_i\). Replay the historical filter and forward equations, and backtrack the matching rows through them to lift each \(b\) to a compatible past assignment. Replay the suffix with a matching row at each next bag and glue both directions. This recovers two actual compatible source assignments on \(X\). A singleton report similarly supplies one assignment of that label and checks that no row in the intersection has a different label. Universal table replay proves the corresponding absence, including initial inconsistency. It does not certify an alleged installed plant's compliance with the factors. That requires the independent source-validity premise.

\textbf{Minimum worst-path audit.}

For an inhabited information state write \(V(z,R)\) for the remaining worst-path cost. Use zero when a singleton report is available. Otherwise evaluate
\[
 V(z,R)=\min_{q\in Q(z)}
    \left\{c(z,q)+
       \max_{a\in\gamma_{z,q}[R\cap\Sigma_{i(z)}]}
                   V\bigl(\delta(z,q,a),R_{\mathrm{next}}(R,q,a)\bigr)
    \right\}
\]
with a minimum over an empty menu equal to \(+\infty\). Each lawful query at a reachable, nonempty state has at least one answer. Public acyclicity orders this recursion even when the current bag stays the same or costs are zero. The proof is reverse induction on public state. At a singleton an immediate common report is proper and cost zero is optimal because all prices are nonnegative. At any other state, a proper uniform audit must pick a query before seeing the reply. A feasible branch must succeed uniformly on the completions of that branch. Its worst remaining cost is therefore at least the optimal subvalue in the formula. These completions are genuine and the maximum over their reply branches is required for the optimum. Conversely select a minimizing query at a finite value and at each feasible reply select a realizing continuation. The separator identity guarantees the same choices work at every historical occurrence of the state. Acyclicity and the finite menus guarantee an attained finite optimum and a proper policy when the formula is finite.

One verified upper certificate is the reported finite policy with each movement and leaf replayed at its \((z,R)\). One matching lower certificate is the complete value table for all information states reachable under legal commands, including infeasible-to-solve states with \(+\infty\). The checker recomputes \(R_0,\Sigma\), success or failure at each stop, every feasible reply and successor, and each exact minimum and maximum. A lower bound cannot be obtained from a value table that omits a lawful candidate or invents an infeasible adverse reply. If the score is infinite, from every admissible query at a nonsingleton state the certificate supplies a feasible response with infinite continuation, until a still-ambiguous state has no commands. This is an adversarial proof of inaccessibility. For any fixed audit the chosen adverse path is realized by at least one fixed source assignment, since every entire prefix is feasible and the path is finite. There is no assumption of changing the installation as replies arrive. A checker that separately validates a feasible finite policy and a lower bound can retain either bound if a larger computation is incomplete.

These exact criteria apply to deterministic trees. Internal randomization cannot reduce a common almost-sure budget for uniform zero-error identification in this finite deterministic-response model. For each compatible installation the set of seeds yielding a correct report within budget has probability one. Their intersection over finitely many installations still has probability one. Fixing a seed in that intersection yields a deterministic audit with the same guarantee.

There are at most \(|Z|2^K\) possible pairs. Generate visited pairs from the initial entry by all legal replies, or use the upper bound and discard uninhabited pairs. Build each update by at most \(k\) successive matching row scans, each on at most \(K^2\) pairs. Factoring a bag uses scanning of its \(K\) rows and its explicit input factors. A hash key or lexicographic bit-vector key on bag rows can store each projected relation, and verified lookup is possible in polynomial time in key length and table address length. The number of fixed register updates and response lengths at a state is bounded by \(L\). Preprocessing, transitions and certificate comparisons use time and size \(|Z|2^K\operatorname{poly}(L,K)\). Price bit lengths grow by at most \(\lceil\log_2(|Z|+1)\rceil\) beyond the largest price bit length since an explicit acyclic path uses fewer than \(|Z|\) steps. Compare the binary requested budget to the exact root value. No numeric audit budget table is needed, and \(K\le d^w\) for \(w=\max |B_i|\) and \(d=\max|\operatorname{dom}(x)|\). The input bound and the fixed-parameter claim include preprocessing and verify time. They do not imply polynomial preprocessing from an arbitrary succinct source to the required factors and public graph, and an operational numeric threshold can enlarge factor domains.

There are arbitrarily large correlated families with bounded \(K\). For example, take \(n\) installed line facts \(x_1,\ldots,x_n\in\{0,1\}\) with no neighboring pair both equal to one, and let the target be parity of the facts. Use auxiliary prefix parity bits \(p_i\), the factors \(p_i=p_{i-1}\oplus x_i\) with \(p_0=0\), adjacent compatibility factors, and a final \(t=p_n\) factor. The bags \(\{t,x_{i-1},x_i,p_{i-1},p_i\}\) with appropriate initial boundary cover all these clauses. Every complete admissible installation determines unique parity auxiliaries. At stage \(i\) the interface can read \(x_i\) or execute a public constant-reply advancement command, with costs and permitted menus carried by a supplied finite acyclic state graph. A physical R-freeze implementation of any fixed Boolean acceptance gate is given below. Correlations admit at least \(2^{\lceil n/2\rceil}\) distinct installations by freely choosing odd positions with even positions zero. Here \(K\le 32\), and the protected audit has a polynomial bound in the explicit linear or polynomial-sized graph input. This is a structural guarantee and no claim of superiority over all variable elimination, decision-diagram or fixed-order solvers.

\textbf{Shared intervention expression and boundary examples.}

Work with occurrence universe \(A=\{(z,q):q\in Q(z)\}\), or a chosen removable subuniverse whose complement is always kept. For each reachable state in the full interface define one value node using the displayed recurrence. Gate each command term by its keep bit, giving \(+\infty\) when disabled. With a singleton stop use a constant zero node. Every outgoing command uses the same already verified reply set and successors as in the unmodified graph. Its deletion is a whole source-authorized effect on availability. The value at the root of this single expression DAG supplies the optimum for any mask by the same induction. Edges and gates are bounded by the expanded exact information-state construction, so the size bound follows. Infinite costs do not introduce a new callable query or branch. Enabling additional commands cannot increase optimal cost, hence the exact budget predicate is monotone in keep bits.

For every proper audit feasible within a specified budget, its support is the set of occurrences used on some feasible path, projected onto the deletable universe when needed. It still succeeds when precisely that support and all undeletable occurrences are kept. This follows by unchanged reply sets and transitions of all used commands. Conversely any proper audit in the deleted interface has support contained in its keep set and is valid in the full interface. Thus minimal sufficient keep sets are the minimal supports. A deletion destroys every proper audit iff it intersects every minimal support. The inclusion-minimal destructive masks give the stated clutter of transversals. If there is no feasible support already, the empty deletion suffices. If the empty support is feasible, no deletion can block closure. A proposed minimal keep set can be checked by one evaluation and one per one-element deletion. Full enumeration of either family can require a large output. These are availability interventions on the \textbf{audit} interface. Resource freezes on physical D-35 tasks are separate and require the whole-atom support semantics used to generate the target.

Three controls delimit the theorem.

\begin{enumerate}
\def\labelenumi{\arabic{enumi}.}

\item
  \textbf{Same projection, different prices.} Initially a public operator can select either of two processing channels through constant-reply, pure commands. Each channel leads to the same bag and posterior relation for an unknown binary target. The same independently recorded instrument is priced at one in one public state and at three in the other. Residual optimal costs differ. Erasing the channel state makes either one estimate invalid. A menu difference gives the same obstruction. A public state claiming to summarize all observable history must include every such dependence, or carry it explicitly in a larger state.
\item
  \textbf{Dynamic inaccessibility without a global pair.} Four installations carry independent \(a,b\), with target \(a\oplus b\). The initial lawful choice is an exact independent read of either \(a\) or \(b\), consuming a single authorization token and leaving no further read. Every opposite-target pair is distinguished by at least one of these two policies, yet whichever one is chosen, both of its reply cells contain opposite labels. There is no uniform identifying policy. Fixed full-menu pair logic cannot treat two mutually exclusive commands as jointly available. The exact value-table obstruction is instead a witness indexed by the legal first choice.
\item
  \textbf{Premises are substantive.} A variable allowed in separated bags without occupying the intermediate separator can be forgotten and then assigned a new incompatible value. For example, a past read of \(x\) and a future read of the same \(x\) must agree on an unchanged register, even if an intermediate gate only refers to another coordinate. The interval requirement keeps the needed connection. A menu depending on unobserved hardware that silently grants a different command is also a different model. Any visible rejection, timing signal or permission observation must be included in the lawful response interface and updates. A destructive probe that mutates source facts needs a new transition model. The fixed latent relation as used here does not apply.
\end{enumerate}

\textbf{M02.2 · Exact cost of correlated lawful routing.}

\textbf{Realization of the fixed operational property.}

A simple canonical construction realizes every gate \(G(\theta)\) specified by an ordinary bounded Boolean circuit on fixed source facts. Let \(X=\{0,1\}\) have two sealed target classes. An admitted history \(h_{\rm open}\) represents both, \(h_{\rm goal}\) represents just \(0\), and two different unused history representatives \(u,v\) have no directly compatible world. At the initially open entry the certificate relation is discrete. A single whole unit-cost dispatch updates \(H\) to \(h_{\rm goal}\). When \(G(\theta)=0\) it also merges the two unused representatives \(u,v\) in \(P_R\); otherwise it leaves \(P_R\) unchanged. Authority is constant for this version. There is one valid deterministic supported result. The authorized objective is still binary closure and the correct atomic R effect compares extensional relations, whether or not the changed class affects a target authorization decision. The unique cheap closer survives R freeze exactly when \(G=1\). An optional always-available E-only price-two backup makes the frozen costs respectively one and two, without changing the threshold. The source command performs one authorized whole dispatch. No separate action exposes either side effect. The \(1\{\kappa^{\neg\{R\}}\le1\}\) value is therefore the gate \(G\).

For M02.2 set \(G=1-g_r\). Fix a source route record and separate site-register records, with at most one blocked site because the deployment has one outstanding site hold. Every route is admitted in each such gate configuration. The routine implements the indexed record test using a finite lookup. The envelope reads take raw site records from independently licensed archival instruments. The route read discloses the signed active site, and the current-site selection enables a specific physical inspection of that site's same gate record. Its price \(s\) is at most the archive-read price. Other cheap inspection commands stay unavailable. Readbacks and choice of inspector change only the later audit menu. The underlying source dispatch and counterfactual property stay fixed.

\textbf{Exact minimax bounds.}

A static sequence must be admissible under every complete installation. It cannot include an inspection of a named site. Before a route read no such command is enabled, and after it there remains another admissible route whose actual queried sequence must be legal. Archive responses do not remove routes from the compatible class, as all routes are allowed for every gate record vector. There is no built-in routing action that implicitly dispatches the inspection query. The route read is indispensable: for any \(j\) use gate vector \(e_j\), with exactly its \(j\)-th bit one, and compare routes \(j\) and \(k\ne j\). Without the route read the same gate records give opposite thresholds. Similarly, omitting archive \(j\) leaves the two installations \((j,e_j)\) and \((j,0)\) with equal responses to the remaining queries but opposite targets. Therefore every static identifying sequence costs at least \(a+\sum_j c_j\). The sequence taking all those readings achieves the bound.

For the adaptive lower bound follow installation \((j,e_j)\) for any \(j\). If it never reads route, it cannot use any inexpensive inspection, and \((k,e_j)\) for \(k\ne j\) gives the same transcript. If it never reads \(g_j\) by either authorized method, \((j,0)\) gives the same transcript, including all dynamic permissions because the route is the same. Both kinds of pair show that this path must acquire the route and the \(j\)-th bit. Its cost is at least \(a+s\). Read route and then query the corresponding cheap inspection to attain \(a+s\) on every installation. These are exact target costs. Their difference and ratio grow unboundedly, for example at \(a=s=1,c_j=2\). The structure combines correlations with genuine query eligibility. Its comparison would change if conditional site execution were an admissible static bundle, if the route were visible initially, or if cross-site gate records determined each target solely from other queries.

\textbf{M02.3 · Matched stochastic acquisition costs.}

\textbf{An operational identity front end.}

Use the same open and closed admissible fibers and discrete root representation as the preceding canonical construction. At the common entry the fast cost-one discharge realizes gate \(x_0\). In addition it alters the extensional authority register in either mode by granting a postreceipt authorization for an independently callable receipt-release action. This changes an actual subsequent permission. It does not affect stopping at the closed checkpoint. Thus the fast path is always removed by A freeze and survives R freeze exactly when \(x_0=1\).

The same entry starts a serial review route with \(d\) forced checkpoints, followed by a price-two E-only closing certifier. At stage \(i\) exactly one of two distinct source review actions is admitted by an ordinary record guard, depending on \(x_i\). Their prices are fixed at \(1\) and \(1+2^i\). A shared extensional authority table and the same input facts persist over the review route. Each stage leaves the authorization target open and advances the raw stage marker, with no resource touch needed. The review certifier changes only \(H\). There are no alternate continuations or earlier closure. All source commands and their individual prices are fixed across installations. Each source transition has one supported successor. Selection of the fast or review route is allowed at the same entry, and the postreceipt command is enabled only after fast closure by the newly issued authority and its ordinary checkpoint guard.

The cheapest review path costs \(2+d+\sum_i2^i x_i\). This is at least three and is the exact A-freeze value. It is injective in all the fact bits and depends on each independently. Under R freeze, budget one succeeds exactly when fast discharge remains available, giving target \(x_0\). The pure audit has no operational plant actions. It has one independently licensed instrument per installed record. All facts have operational significance at the same entry even when ignored by the tested threshold.

With noiseless certain answers, one read suffices and is necessary for \(x_0\). Every certain audit of the whole vector must read every slot on every path, since an unqueried bit can flip without changing any readback on that path. Reading all bits attains this bound, and adaptive access cannot improve it. Choosing to ask only the primary instrument says nothing about the external costs of the review route.

\textbf{Noisy upper bounds and posterior factorization.}

Write \(\Delta=1/2-\eta>0\) and \(\lambda=\log((1-\eta)/\eta)>0\). For a task involving \(s\) bits, where \(s=1\) for the target and \(s=d\) for full identity, take \(m=\lceil\log(s/\delta)/(2\Delta^2)\rceil\) samples from each of the needed instruments. A majority estimate of a fixed bit fails with probability at most \(\exp(-2m\Delta^2)\) by the bounded-sum Hoeffding inequality. The union bound gives at most \(\delta\) probability of any error on the whole report, for every installation. This uses \(sm\) calls with a deterministic path cap and the same expected cost.

For the lower bounds equip the admissible vectors with the uniform product prior. Uniform correctness implies Bayes average error at most \(\delta\) under this prior. Observe the independent auditor seed in the proof in addition to the audit history. Conditional on any fixed seed, the chosen instruments, the decision to stop, and its report are functions of the response sequence and the protocol. For a finite complete history, the likelihood of that history at a latent vector \(x\) is either zero from a mismatch with the fixed policy or is the product of the likelihoods of each recorded BSC response at its queried coordinate. The mismatch condition is independent of \(x\) for the same history, and an actual observed finite transcript has positive likelihood under every \(x\). Repeated calls, adaptive choices and the finite stopping record have no further latent factor once the seed and complete transcript are fixed. This establishes factorization of the posterior across bits at any realized finite stop. For arbitrary independent seeds the calculation applies for almost every seed and then integrates. Almost-sure finite stopping under every completion removes any need for a conditional formula at an infinite transcript.

Let \(n_i\) be the number of calls at bit \(i\) in the full transcript and let \(N=\sum_i n_i\). Its posterior log odds has absolute value at most \(\lambda n_i\), since each observation changes the odds by one factor \(\exp(\lambda)\) or its inverse. Write \(e_i\) for the smaller posterior probability of the two bit values. Then
\[
 e_i\ge \frac{1}{1+\exp(\lambda n_i)}\ge \frac12\exp(-\lambda n_i).
\]
For full identity the best posterior success probability is the product of the coordinate maxima. It is at most \(\exp(-\sum_i e_i)\), using \(1-u\le \exp(-u)\), and convexity of \(x\mapsto \exp(-\lambda x)\) gives
\[
 \sum_i e_i \ \ge\ \frac{d}{2}\exp(-\lambda N/d).
\]
For the target bit alone the same bound with \(s=1\) holds since \(n_0\le N\). Thus every report, conditional on the full transcript and seed, has posterior error at least
\[
 1-\exp\!\left(-\frac{s}{2}\exp(-\lambda N/s)\right).
\]
This lower bound is decreasing in \(N\). It applies to any adaptive sequence and any terminal report under the declared model.

\textbf{Worst-path versus worst expected cost.}

If every completion has \(N\le B\) almost surely, then this conditional lower bound and the Bayes requirement give
\[
 B\ \ge\ \frac{s}{\lambda}\log\frac{s}{2[-\log(1-\delta)]}.
\]
For the worst expected criterion take \(M=\max_x \mathbf E_x N<\infty\). If \(M=0\), no calls occur under the prior and the success guarantee already fails for \(\delta\le1/8\). Otherwise the Bayes expected number of calls is at most \(M\), so Markov's inequality gives \(\Pr(N\le2M)\ge1/2\). On that event use the same conditional error lower bound at \(2M\) and integrate. Then
\[
 \delta\ \ge\ \frac12\left[1-\exp\!\left(-\frac{s}{2}\exp(-2\lambda M/s)\right)\right],\qquad
 M\ \ge\ \frac{s}{2\lambda}\log\frac{s}{2[-\log(1-2\delta)]}.
\]
For \(0<\delta\le1/8\), \(2[-\log(1-2\delta)]\le (16/3)\delta\), and \(\log(s/\delta)\ge\log 8\). Subtracting the fixed constant \(\log(16/3)<\log 8\) leaves a uniform positive multiple of \(\log(s/\delta)\). Both criteria therefore have the claimed exact asymptotic orders with constants depending only on \(\eta\). The upper bound uses the very same menu, prices and uniform confidence constraint. It does not assume an average prior guarantee as the desired criterion.

If the replies share one permanent corrupted readback coin per instrument, repeated measurement need not reduce error at all. In a single-bit example the first read has BSC error \(\eta\), every later read equals it, and a uniform guarantee for \(\delta<\eta\) fails by the two-point Bayes lower bound. This violates the independence premise. At positive independent noise, every finite stopped history of positive probability under one fixed vector has positive likelihood under every vector with the same policy seed. Zero-error uniform identification with a.s. finite stopping is therefore impossible for a nonconstant target. These boundaries explain why the deterministic exact staged solver does not imply finite certain stochastic identification.

\textbf{M02.4 · Two representations and their complexity.}

\textbf{Opposing-pair existence.}

For a Boolean-circuit compatible predicate and exact Boolean-circuit target, a certificate for existence of an opposite pair supplies two polynomial-length latent vectors. Evaluate the compatibility circuit on both and verify their target disagreement in polynomial time. With no queries this is exactly nonconstancy on the compatible domain. For hardness take any 3-SAT formula \(\varphi(y)\) and an independent switch \(b\). Permit all assignments to \((y,b)\) and use target \(b\land \varphi(y)\). The compatible witness \(b=0,y=0\) is explicit, and a target-one assignment exists iff \(\varphi\) is satisfiable. The gate realization from M02.2 compiles this Boolean expression into a fixed-atomic typed one-step source with the target the same-entry R-freeze threshold, using a polynomial-size description. This is a polynomial reduction. Complementing gives coNP-completeness of zero-query point identification on the known nonempty representation. This tests existence of a certified disagreement independently of the cost of synthesizing any acquisition.

\textbf{Exact dynamic hardness with succinct public registers.}

For membership in PSPACE, recursively search candidate legal actions for a winning audit of remaining budget. Store the current explicit transcript and public registers along the depth-\(H\) recursion. At a feasible history decide whether there is a common target by enumerating the polynomial-bit latent vectors, checking compatibility and the entire sequence of deterministic replies. Enumerate candidate replies to any lawful command using their polynomial-bit encoding. A reply is feasible if a single latent vector satisfies all past answers and this new one. Each test uses polynomial workspace by enumerating one vector at a time. The encoded price bound and remaining budget need polynomial bit lengths and the recursion depth is at most the unary \(H\). At each nonsingleton history with remaining depth, accept if some permitted command costing at most the remaining budget has all feasible reply branches accepted with that price charged. A singleton can stop immediately. Induction on remaining depth proves the search exact. The exponential time and large number of prefixes do not increase space beyond the polynomial stack. Input inconsistency can be detected independently by the same latent scan if not promised.

For hardness start with the quantified Boolean formula
\[
 Q_1u_1\cdots Q_m u_m\,\varphi(u_1,\ldots,u_m).
\]
For each universally quantified bit introduce one independent latent record, and introduce an additional independent target bit \(b\). The compatible class is the entire product. Design a fixed acyclic public interface with counter and \(m\) public audit registers initially blank. At stage \(i\) the allowed commands depend only on the prefix quantifier:

\begin{itemize}

\item
  At an existential quantifier, allow a public write of either Boolean choice by two distinct unit-price commands with a constant reply and with this choice recorded solely in the \textbf{audit} register. They do not change the fixed installation, the universal latent data or the target. These can be requests selecting either of two testing profiles.
\item
  At a universal quantifier, the only command is the exact unit-price measurement of that stage's independent latent record. Store its reply in the public register. Both replies are feasible at every prefix through that stage.
\end{itemize}

Each of these commands advances the stage by one. After the \(m\)-th choice the sole unit-price instrument for measuring \(b\) is callable iff \(\varphi\) evaluates true on the public record. This is a supplied policy gate encoded by the formula's Boolean circuit. There are no shortcut queries, and no earlier response depends on \(b\). The target uses the fixed-gate canonical one-step R-freeze plant with gate \(b\). Thus every feasible transcript before the final measuring command still permits both target values even if it determines all other latent facts. Use \(H=B=m+1\). The transition and permission system is representable by polynomially many bits and polynomial-size circuits. Hidden universal values are picked once as part of the underlying completion. The public eligibility formula uses observed universal values and the tester's own past choices only.

If the quantified formula is true, its existential strategy selects the public profile commands from the previous observations. It satisfies the final gate on every admissible response sequence and then reads \(b\) within budget. Conversely any identifying policy must eventually use the \(b\) instrument on every installation since no other reply separates its values. The choices at existential stages depend only on the earlier replies. Thus those choices supply an existential strategy winning against every universal record vector. Running through independent universal bits gives exactly the QBF game in the original order. The formula is true. This proves PSPACE-hardness and hence completeness for the stated representation. It relies on succinctly specifying the permission gate and public registers. Explicitly listing their possible values can require exponential space, so it does not contradict M02.1. It also leaves the accepted T10 binary SUBSET SUM reduction intact for a different paired numerical oracle.

\textbf{M02.5 · Checked approximation and a succinct representation boundary.}

\textbf{Certified partial expansion.}

Retain the lawful query and feasible answer checks at every expanded state. At a frontier state the exact value is nonnegative and at most \(+\infty\). Any independently certified bounds may tighten this interval. Addition of a fixed nonnegative price, maximum over the feasible replies, and minimum over lawful queries and any allowed stop are all nondecreasing in every continuation value on \([0,+\infty]\). Reverse induction in the expanded graph therefore gives a lower bound by substituting the lower endpoints and an upper bound by substituting the upper endpoints. Unexpanded states require no claim of optimality. Validate the frontier annotations or use the trivial interval. A witness for a finite upper value consists of a fixed chosen query, a witness continuation at every feasible reply, and at any unexpanded state reached by that witness a proper bounded-cost audit with its own checks. Thus the finite upper value can be achieved or improved by a verified complete audit if the stated upper witnesses are supplied. If \(U\) is such an upper budget and \(L>0\) is a checked lower bound, then \(U/\mathrm{OPT}\le U/L\).

At every state intersect its previous verified interval with any expanded recurrence interval. Extrema and intersections preserve the exact value, and graph ancestors inherit only tighter bounds since the expression is coordinatewise monotone. Store the old realizing upper witness if its bound is used after an intersection. Without this retention, expansion of a previously bounded unknown state with all of its children initialized to \([0,+\infty]\) can loosen a known bound. For example a sole legal first command costs one and leads to a nonsingleton child whose only useful query costs two. Independent verification gives the original interval \([3,3]\). Expanding only the first state and forgetting that bound gives \([1,+\infty]\). A partial table gives no promised finite multiplicative ratio.

\textbf{Inapproximability with a finite fallback.}

Take the QBF construction used for M02.4 and add a named independent read of the same installed \(b\)-record at the initial and every preterminal public state where a query can still fit in the same horizon. Charge it \(D=2^n\), where \(n\) is the input length of the quantified formula. The direct read ends the audit. It is part of the supplied permitted menu, independent from the ordinarily priced read available after eligibility is satisfied. It has the same target-independent physical measurement interpretation. Query and register circuits remain polynomial-sized since the binary encoding of \(D\) uses only \(n+1\) bits. The unchanged unary horizon \(m+1\) permits the fallback immediately.

On a true QBF, the unit-price policy costs \(m+1\), and the only alternatives reading the record have cost at least that value for sufficiently large \(n\). On a false QBF, every proper audit uses the fallback on at least one feasible path. Otherwise wholly avoiding it would give the previously disproved strategy. Thus the worst cost is at least \(D\), and the initial fallback attains exactly \(D\). Both optimum values are positive finite.

Suppose an estimator provides \(\mathrm{OPT}\le\hat V\le p(L)\mathrm{OPT}\) in polynomial time for some fixed polynomial \(p\). The constructed input bit length \(L\) is polynomial in \(n\). Hence for all sufficiently large \(n\) it satisfies \(p(L)(m+1)<2^n\). The estimator decides QBF by testing \(\hat V<D\), and the finitely many smaller cases can be handled by an exact finite procedure without changing polynomial time. Since QBF is PSPACE-complete, this yields \(\mathrm P=\mathrm{PSPACE}\). The argument preserves deterministic measurements, a fixed operational threshold, a known finite option and a polynomial horizon. It relies on binary prices and succinct public permission states. It gives no impossibility of obtaining instance-specific intervals or of stronger approximations under additional structural promises.

\textbf{Scope and checks.}

The optimizer needs the complete operative answer record, the full compatibility constraints and executable commands. Physical closure, stochastic risk under an unknown law, and arbitrary source-to-factor extraction are separate mathematical tasks. Neither the \(K\)-dependent upper bound nor an independently reconstructible small check proves tractability of the question with general unbounded paths, target or model misspecification, arbitrary setup cost changes or combinatorial overlap between operational actions and evidence acquisition. More efficient width regimes, the best structural approximation guarantees and useful state refinements for stochastic dynamic audits remain open.

\end{proof}

\subsection{M03. Joint risk and sufficient anchor facts}

\label{ext:m03}

Marginal calibration intervals can discard restrictions that connect one stage or policy to another. Jointly feasible path masses enforce those restrictions before an extremum is taken.

\begin{theorem}[M03.1. Joint trajectory risk from local conditional and event restrictions]\label{ext:m031}

Fix finitely many runs \(d\), contexts \(x\), outcome letters \(y\), event labels \(b\in\{0,1\}\) and installation modes \(m\). Every run fixes a lawful controller, initial distribution assumptions, finite horizon, admissible support and a trusted event interpretation. Within a mode let \(q_{dxyb}\) be deployment masses, \(a_{dxy}=\sum_bq_{dxyb}\) and \(w_{dx}=\sum_ya_{dxy}\). Reference masses \(s_{dxy}\) satisfy \(\sum_y s_{dxy}=w_{dx}\), and \(\sum_x w_{dx}=1\). All masses are nonnegative. For each occupied context the reference conditional row satisfies fixed affine constraints \(H(s/w)\le h\) and the two conditional rows are within total variation \(\delta_{dx}\). Occupied outcome cells have checked conditional event probability in \([l_{dxy},u_{dxy}]\). Additional joint restrictions are closed affine constraints on these arrays and bounded auxiliary variables. All named policy flows, data snapshots and cross-run restrictions must be imposed jointly.

Each mode is represented by a compact polytope. Put \(r_d=\sum_{x,y}q_{dxy1}\) and \(R=\max_dr_d\). Its image on each inhabited mode is a closed interval \([L_m,U_m]\). One epigraph linear program computes \(L_m\); maximizing each \(r_d\) and taking the maximum gives \(U_m\). The full image is the union of the intervals. At threshold \(\tau\), uniform passing requires \(\max_mU_m\le\tau\) and uniform failure requires \(\min_mL_m>\tau\). A mixed result has two witnesses. No inhabited modes means inconsistent evidence. Operational sharpness requires every feasible array to be realizable; a verified outer inclusion alone supplies sound bounds.

For \(K\) prespecified faithful Bernoulli calibration cells, let \(n_i\) independent draws with the designated mean estimate cell \(i\). Simultaneous radii \(\rho_i=\sqrt{\log(2K/\alpha)/(2n_i)}\) suffice for \(0<\alpha<1\), even for dependent batches across cells. A checked bridge \(|p_D-p_E|\le\eta\) permits expansion by \(\eta\) and clipping to the unit interval. Selecting registered constraints preserves actual membership with probability at least \(1-\alpha\), conditional on valid source, policy, support and scoring premises. Homogeneous issuance then has unconditional wrong-issuance probability at most \(\alpha\). Exact numerical optimization uses rational input or certified outward rounding.

\textbf{Conditions and boundary cases.} Reference restrictions use an explicit finite, fixed affine matrix for each declared run, policy and data snapshot. Exact LP computations use encoded rational coefficients. A realized joint-source family needs compatible flows linked by its actual reconstruction constraints. Arbitrary selectors, gaps between compatible row choices, or independently optimized per-policy laws can change sharpness. A confidence set must contain the true linked model under its stated assumptions.

\end{theorem}

\begin{theorem}[M03.2. Minimum sufficient independently checked anchor facts]\label{ext:m032}

Give each of \(k\) eligible pure anchor reads a nonnegative acquisition-plus-verification price \(c_i\) and mode-specific closed affine fact \(F_i\) in a jointly consistent realized record. A subset \(J\) certifies the proposed label exactly when it meets \(D(z)=\{i:z\notin F_i\}\) for every opposite-label completion \(z\). Minimum cost is the weighted transversal optimum, with infinity for an empty edge. At most \(2^k\) different patterns need be generated. For a prospective audit, specify all jointly possible full answer tuples, with compatible arrays in each mode forming closed polyhedra and projecting onto every admitted completion. The array and its tuple stay fixed during the audit. In a nonempty homogeneous state the value is zero. Else minimize \(c_i\) plus the largest continuation value on a feasible answer branch over the unread items; infeasible resolution costs infinity. The minimum is attained when finite. A faithful direct read available everywhere at constant cost \(d\) changes an exact semantic-only optimum \(W\) to \(\min(d,W)\).

\end{theorem}

\begin{proof}

At a fixed run and context, \(a,s\) have the same total \(w\). Multiplying the fixed affine reference restriction by \(w\) gives \(Hs\le hw\). Introduce \(0\le t_y\le1\) and require
\[
-t_y\le a_y-s_y\le t_y,\qquad \sum_y t_y\le2\delta w,\qquad
l_ya_y\le q_{y1}\le u_ya_y.
\]
At positive occupancy division gives the original conditions, while \(t_y=|a_y-s_y|\) proves the converse. At zero occupancy all those masses vanish and no conditional row is inferred. The constraints hold simultaneously across the roster; retain all original affine links, rather than optimizing contexts independently. One compact polytope results per mode. A finite maximum of linear functionals is continuous. Its extremizers exist and the segment between them, by continuity, achieves every intermediate value. Minimizing an epigraph variable constrained above every \(r_d\) and maximizing each coordinate separately yield the claimed endpoints. Feasible primal solutions are physical witnesses exactly under the stated realization premise.

The Bernoulli tail bound makes each failure probability at most \(2\exp(-2n_i\rho_i^2)=\alpha/K\). On the intersection of all \(K\) good events, every selected fact from the registered family contains the same true completion; the bridge expansion is the triangle inequality. Batch choices that change the row being measured would need fresh justification.

For a proposed label, an opposing array survives precisely when it satisfies every selected fact. This proves the transversal condition. Solve for a cheapest transversal of the edges seen so far. Test its cells for an opposite array: impose \(r_d\le\tau\) for all \(d\) to test passing and maximize individual risks to test strict failure. A feasible opposite array returns an unhit edge disjoint from this transversal, necessarily new. An empty edge rules out every certificate. Otherwise insert the edge and repeat. A transversal surviving this test is feasible for the full problem and already minimizes over a relaxation, hence is optimal. The number of possible edges bounds the number of insertions. Conversely an arbitrary finite weighted hitting set embeds by one always-safe mode and one unsafe mode per edge, with a realized zero read at index \(i\) contradicted exactly in modes whose edge contains \(i\).

For the prospective problem, filter the entire tuple-and-array relation by the actual partial answers. This preserves correlations even when individual fact intersections would be larger. Induction on unread items gives the displayed recursion: any first read must pay its price and succeed on every viable child, while combining one optimum subtree per response attains that maximum. A branch has a single joint witness by construction. There are at most \(\prod_i(1+|A_i|)\) partial tuples for answer alphabets \(A_i\). Any tree using the direct read on a reachable branch costs at least \(d\) in the worst case; one never using it costs at least \(W\). Taking the corresponding root alternative achieves the lesser.

\end{proof}

\subsection{M04. Audits with strategic responses}

\label{ext:m04}

A robust audit is a game of imperfect information. Its private random choices have value only when the responder must commit to a legal strategy without seeing them.

\begin{theorem}[M04.1. Finite exact audit optimization]\label{ext:m041}

Fix a finite authorized audit tree of horizon \(H\) with perfect recall, finite actions, specified information sets, chance rows and all actionable records. A pure responder plan chooses a model with fixed deployment label before initialization, can depend on its own information, and cannot inspect a private future auditor draw. Any shared or disclosed random signal is represented explicitly. The auditor has common callable commands at its information sets. Declare nonnegative path costs and terminal losses, including wrong verdict or abstention for the applicable fixed class. Against every allowed responder impose each class-specific failure cap, expected-cost cap, or zero chance of exceeding a hard path budget. Mixtures over the two players' pure plans use independent private randomization.

Enumerating pure plans gives an exact linear program for feasibility, optimal worst expected cost or minimum worst error. For rational data a finite optimum has rational witnesses; infeasibility has a rational certificate. More compactly, if the permitted strategies are exactly the behavioral strategies of the given perfect-recall tree and their private mixtures, there is a formulation polynomial in its explicit rational encoding. Let \(Ex=e\), \(Fy=f\) and \(x,y\ge0\) be the two sequence-flow systems. For each chance-weighted terminal loss matrix \(A_j\) and cap \(b_j\) require a free vector \(v_j\) with \(F^\top v_j\ge A_j^\top x\) and \(f^\top v_j\le b_j\). The same auditor vector \(x\) serves every loss. Coupled responder restrictions require a representation that preserves them.

\textbf{Conditions and boundary cases.} Each pure responder alternative fixes the true class before the initial chance record. Mixed drawing of pure plans is a minimax device. The chosen pure alternative cannot anticipate the private auditor seed. Conditioning a mixed-label initial law requires class-normalized or centered constraints. An anticipating adversary can erase the signal in the protected-bit example.

\end{theorem}

\begin{theorem}[M04.2. Exact composite obstruction for a fixed audit]\label{ext:m042}

Fix an audit design through its last observation, before its verdict. Let \(T\) be its finite complete record alphabet and \(\mathcal P,\mathcal Q\) the nonempty compact convex hulls of its two class laws. A randomized verdict with uniform class failures at most \(\alpha,\beta\) exists exactly when, for every \(\lambda\in[0,1],P\in\mathcal P,Q\in\mathcal Q\),
\[
\sum_{t\in T}\min\{\lambda P(t),(1-\lambda)Q(t)\}
\le\lambda\alpha+(1-\lambda)\beta .
\]
Strict excess witnesses infeasibility. On finite rational hulls its largest possible excess is the optimum of a linear program. The criterion also applies if abstention counts as failure in either class. Intersecting hulls force maximum class failure at least \(1/2\); disjointness alone need not meet prescribed caps.

\end{theorem}

\begin{theorem}[M04.3. Interval responders and stopping-valid confidence]\label{ext:m043}

In the interval model of Theorem~\ref{thm:t24}, its separation condition and explicit sample quotas characterize finite-cost uniform error below any prescribed \(0<\delta<1/2\), under its full-record richness assumption for necessity. In a finite capped score-only tree with unrestricted history-wise interval means, endpoint choices followed by Bernoulli chance draws realize all the permitted laws, so M04.1 computes the optimum. For a fixed design on finite record alphabets at each depth, let \(M\) be any consistent reference path law. Retain an admitted law \(P\) along a prefix \(h\) only when \(P(h)>0\) and \(M(h)/P(h)<1/\delta\). The probability that the true law is ever excluded is at most \(\delta\) under each admitted actual law, also for a composite family. Homogeneous issuance controls unconditional wrong issuance. With finite extreme laws the retained-class feasibility checks are linear inequalities in law weights with strict inequalities decided by positive common slack.

\end{theorem}

\begin{theorem}[M04.4. Expected-cost horizon approximation]\label{ext:m044}

Fix a finitely branching infinite audit with perfect recall, legal nonanticipating behavioral rules, fixed class labels and consistent truncations. Each finite responder plan extends to an allowed infinite one and conversely each infinite plan restricts to an allowed finite one. Every continuation completes another probe at price at least \(c_0>0\), stopping costs zero, and every forced declaration at a cap is legal in the original game without changing its preceding cues. If an almost surely terminating policy meets worst expected cost \(B<\infty\) and class failure caps \(\alpha,\beta\), truncation at \(H\) probes meets the same cost and increases each failure cap by at most \(B/[c_0(H+1)]\). Conversely, if for every \(\epsilon>0\) a finite cap achieves \(B,\alpha+\epsilon,\beta+\epsilon\), a single almost surely terminating infinite policy achieves the exact three bounds. No finite deterministic stopping time at the limiting caps is implied.

\textbf{Conditions and boundary cases.} For exact finite-cap comparison and the limiting existence argument, cap declarations must be actual legal terminal moves on the corresponding histories of the original model. Their failure cost includes wrong verdicts and abstention according to the declared loss. The exact limiting converse also needs a properness condition, supplied here by the positive lower bound on every probe price. Granting a special stop action only inside a truncated problem can make every cap feasible without creating any such infinite-model policy.

\end{theorem}

\begin{proof}

Index pure auditor and responder plans by \(i,j\). For each terminal loss \(L\) let \(M^L_{ij}\) be its expectation under the two plans and the fixed chance nodes. Private contingent random choices are distributions on pure plans: sample the private choices before play and fill off-path instructions arbitrarily. Conversely a private seed implements any such mixture. For auditor weights \(p_i\), each row \(\sum_i p_iM^L_{ij}\le b_L\) is necessary because \(j\) is admissible, and together these rows are sufficient by linearity. Each pure responder fixes its class before initialization, so checking its class loss uses the proper denominator. Zero mean of a nonnegative hard-budget violation indicator excludes every positive-probability overrun. The simplex is compact and the constraints closed, proving attainment for a feasible finite optimization.

For the smaller formulation, perfect recall gives each information set \(I\) one preceding own-action sequence \(s(I)\). A realization vector assigns unit mass to the empty sequence and obeys \(\sum_a x_{Ia}=x_{s(I)}\), \(x\ge0\). A pure-plan mixture gives these equalities by disjoint choice at \(I\); conversely use \(x_{Ia}/x_{s(I)}\) whenever the denominator is positive and any legal rule otherwise. Induction recovers every sequence weight. The probability of a leaf is its chance factor times the two last sequence weights. Thus worst loss \(j\) for fixed \(x\) is \(\max_{Fy=f,y\ge0}x^\top A_jy\). Linear programming duality gives exactly the constraints on \(v_j\) in M04.1. There is one row per information set and one coordinate per information-set action, so the explicit finite tree controls the size.

For the fixed-design claim let \(\phi(t)\) be the probability of declaring the positive class. Its errors are \(P(1-\phi)\) and \(Q\phi\). Replacing an abstention by a fixed binary verdict cannot worsen either error when abstention is failure. Put \(u=\lambda P\), \(v=(1-\lambda)Q\). These weighted pairs form a compact convex set: normalize a convex combination separately in each nonzero class. The largest of the two class-cap excesses for fixed \(\phi\) is the maximum, over such pairs, of
\[
\sum_t [u(t)(1-\phi(t))+v(t)\phi(t)]-\alpha\sum_tu(t)-\beta\sum_tv(t).
\]
This expression is separately affine on compact convex domains; finite-dimensional minimax exchanges the maximum and minimum. Now each record chooses the lesser of \(u(t),v(t)\), giving M04.2. On finite lists use nonnegative mixing weights of combined mass one and maximize the sum of overlap variables \(0\le o_t\le u(t),v(t)\) less the weighted caps. A shared law at \(\lambda=1/2\) gives \(\alpha+\beta\ge1\).

For the interval result, the concentration and same-law arguments in Theorems~\ref{thm:t24} and~\ref{thm:t15} apply unchanged to the adversary's actual recorded histories. Random mixtures of the two interval endpoints produce every Bernoulli mean between them; a conditional mixture stays within that interval, proving both directions of the finite tree representation. For an admitted true law \(P\), define \(Z(h)=M(h)/P(h)\) on prefixes of positive \(P\)-mass. Summing over the next positive children gives \(\mathbb E_P[Z_{t+1}\mid h]\le Z(h)\) since the reference may place mass on discarded children. The initial mean is at most one. The first-crossing bound for this nonnegative supermartingale gives \(P(\sup_t Z_t\ge1/\delta)\le\delta\). In an extreme-law simplex each tested prefix imposes the strict affine inequality \(P(h)>\delta M(h)\); finitely many strict constraints hold iff their maximum common slack is positive.

Finally couple the original rule and its truncation through their first \(H\) probes against any allowed responder. Their failures differ only on paths taking at least \(H+1\) probes, whose original cost is at least \(c_0(H+1)\). Markov's inequality gives the forward error slack, and truncation cannot increase cost. For the converse take legal finite rules with errors tending to the two desired caps. Extend each to the same infinite tree, then choose one subsequence converging coordinatewise at every information set by compactness of each finite simplex and diagonalization over countably many sets. Against any \emph{fixed} responder, every finite prefix probability and finite accumulated expected cost is continuous in the finitely many involved auditor coordinates. The common limiting policy therefore meets each finite-prefix bound against that responder. Monotone convergence gives the total expected cost bound \(B\). Infinite continuation would incur infinite cost, so it has probability zero. Increasing finite terminal failure events then give the two exact error bounds. The subsequence was selected from the auditor policies before fixing the responder, establishing uniformity.

\end{proof}

\subsection{M05. Joint monitoring and event fidelity}

\label{ext:m05}

A monitor output is observable; the physical event can remain hidden. Joint event constraints, native transition rows and linked policy semantics answer three different questions about that gap.

The input lists finitely many mutually exclusive semantic modes. A mode can specify several populations or contexts as separate trees. Fix a finite explicit rooted tree for each population in each mode. Its nodes are full typed physical prefixes, including all control memory needed to specify conditional rows, and its edges are whole source moves and their supported outcomes. The synthetic root may have one initial sampling move. Declare authorized controller observations, source actions, admitted history, representation, authority, freezes and the stop and padding convention. Each fixed behavioral policy gives a distribution on callable source actions from its admitted observation history and its declared controller state. At each allowed concrete prefix, check authority and freeze touch for each action receiving positive policy probability over all successor outcomes declared possible at that prefix. The specified menus must be valid uniformly for the allowed rows of that mode. The input can use separate modes when those source premises differ. Path construction retains complete successors under the supplied source contract. Physical stopping and target scoring have the stated finite terminal semantics.

Each nonterminal prefix \(h\) has an allowed rational polytope \(Q_h\) of distributions on its outgoing edges. It includes the fixed action marginals at \(h\), any applicable trusted response-row cuts, and support restrictions. An empty allowed row is permitted at a prefix of mass zero. Each terminal carries a fixed trusted truth vector scored from its physical history, a monitor value determined from its visible history, and any recorded audit or label-availability symbols. The outcome law, including which truth-bearing history occurs, can be unknown. The source-defined target and the independent validation of the evidence constraints are premises. Let \(x_h\) denote node mass and \(f_{h,e}\) denote the mass of edge \(e\). Impose root mass one per tree, \(x_{\mathrm{child}(e)}=f_{h,e}\), nonnegativity, and \(\sum_e f_{h,e}=x_h\). For every row constraint \(A_hq\le b_h\) use \(A_h f_h\le b_h x_h\), with equalities handled in both directions. In each mode also impose the supplied rational affine constraints on masses, including constraints across contexts when independently justified.

An admissible completion selects one mode and one executable law per tree in that mode. At each positive prefix its conditional row is in \(Q_h\). The laws satisfy the supplied affine evidence within that mode. Their scalar weights can be arbitrary real probabilities. Mode choice is fixed across its populations. Randomizing the selection of modes is allowed only if separately represented by an admitted source model. The claimed sharpness is relative to the completeness of this declared input for the intended class. For a target threshold, a consistent class passes when every completion is at or below the threshold, fails when every completion is strictly above, and is unresolved otherwise.

\begin{theorem}[T25.J. Joint monitor identification]\label{ext:t25j}

For every input under the common model, the union of its mode flow polytopes is exactly the collection of feasible tuples of population laws described by their prefix masses. For every finite vector of affine statistics of the masses, its image over this union is the sharp joint identified set. Every point has a realizing law in one admitted mode. In particular a population can report actual violation mass \(\Pr(V=1)\), missed mass \(\Pr(V=1,M=0)\), false-alarm mass \(\Pr(V=0,M=1)\) and any cross-context affine statistic without multiplying marginal feasible sets. Truth and alarm are the separately specified terminal functions. Exact labelled submasses and visibility-only masses can be enforced together.

For any event-conditioned statistic with numerator an affine function \(n\) and denominator an affine event mass \(d\), assume \(n=0\) wherever \(d=0\) on the feasible set. On positive-denominator completions the ratio \(n/d\) has a sharp set that is a finite union of closed intervals, one for each mode with some positive denominator. Endpoints are attained by admissible laws. Undefined completions exist exactly when some feasible mode attains \(d=0\). This applies to \(\Pr(M=0\mid V=1)\) with \(n=\Pr(V=1,M=0)\), \(d=\Pr(V=1)\). In a consistent class with a positive-denominator completion, a universal upper cap \(c\) on the defined ratio holds exactly when every feasible mode has maximum \(n-cd\le 0\). Whether zero denominators remain possible is reported separately.

Mode feasibility, extrema of affine statistics, extrema of such ratios and extrema of worst-context risk \(\max_{d\in\mathcal D}\Pr_d(V=1)\) over finitely many supplied contexts can be found with polynomially many rational linear programs of polynomial explicit input size. The lower worst-context endpoint uses the same completion across all contexts. Extreme affine values and violating caps have feasible law witnesses. Empty feasible union returns inconsistency. No compact description of the projected polytope by all its facets is promised.

\textbf{Conditions and boundary cases.} Numerical LP optima use rational encoded objectives and caps, together with the declared linear model. The assertion of exact event containment is about path outcomes positively realizable under a feasible law. An ambient edge or event pattern can be syntactic while every jointly feasible law gives it zero mass. Where \(P(V=1)=0\), the conditional false-negative rate is undefined.

\end{theorem}

\begin{theorem}[T26.C. Event fidelity under history-conditional coupling]\label{ext:t26c}

Let two native finite physical policy trees have exact initial and complete-prefix conditional rows, incorporating each arm's own fixed legal policy. Each may have a different policy and admitted observation interface. They share a finite comparison horizon after semantically valid padding. Form a product tree of paired prefixes. At each product prefix require that each outgoing row projects to the supplied native row conditional on that arm's own prefix. This includes the initial row at the synthetic root. The coupling row can depend on both prefixes. Intersect the row polytope with any independently justified paired support or affine row restrictions. If no such restriction is supplied then all couplings of the two rows are allowed. A terminal can carry any fixed bounded rational cost of the two true terminal histories.

For every such pair of native laws and every collection of these product row restrictions, the product flow is feasible exactly when a coupling of the two complete laws exists with the stated conditional rows and restrictions at every positive paired prefix. Its optimum and attaining coupling for any terminal cost, or infeasibility, can be computed by a backward recurrence with one local feasibility/optimization problem per product prefix. A feasible row at a prefix must give zero mass to children admitting no feasible continuation. The recurrence first tests feasibility under those required zeros, then optimizes the expected child cost using that row polytope. The root value is the required optimum. This recurrence presumes only the stated per-prefix restrictions. Intertemporal or cross-population affine constraints are handled by the common flow problem.

Take independently scored binary events \(V_D,V_P\) for the original and probed policy runs. On any feasible paired law, writing \(b^+=\Pr(V_D=1,V_P=0)\) and \(b^-=\Pr(V_D=0,V_P=1)\), the signed change from probed to original risk is \(b^+-b^-\). The minimum disagreement cost for \(1[V_D\ne V_P]\) is zero exactly when the specified row-preserving class admits an event-agreeing coupling. Its minimum gives an existential bound on the absolute difference of native risks and its maximum bounds disagreement uniformly over the paired class. Here a trusted score of the probed event and preservation relative to the named original policy are separate premises or checks.

At every positive paired prefix, each arm's conditional continuation law projected onto its native future equals its native continuation law starting from that arm's prefix. Consequently the same signed change and disagreement bounds apply to future terminal outcomes given that paired prefix. For a cylinder \(B\) formed from a union of disjoint paired-prefix cylinders, the conditional native law on \(B\) is the occupancy-weighted mixture of the corresponding native prefix continuations. The common flow yields sharp conditional discrepancy ranges, with the denominator checks of T25.J. This conditional conclusion concerns the declared paired class.

\end{theorem}

\begin{theorem}[T27.J. Joint event and local-risk identification]\label{ext:t27j}

In the common model let a tree describe the actual interacting physical plant at its fixed legal policy, with independently scored terminal global event \(G\) and local events \(L_1,\ldots,L_k\). The full prefix can retain shared resources, communication, authority changes and component memory. Optional paired coordinates for independently audited isolated reference laws \(I_i\) must satisfy T26.C's conditional projection equations for those reference coordinates. An uncertain actual interacting coordinate may have the declared convex row constraints at a paired prefix. Its fixed policy and legal response-row requirements at its own prefix must hold for every admissible paired row there. Rows on the isolated coordinates are checked against their own reference prefixes regardless of the actual plant's state. The product can carry independently supported common tickets or other paired row restrictions. The conditional law of any isolated coordinate on a prefix cylinder is as in T26.C.

For every such input, the joint identified set of all \((G,L_1,\ldots,L_k)\) pattern masses, the uncovered global-event mass \(u=\Pr(G=1,L_1=\cdots=L_k=0)\), the global risk \(r=\Pr(G=1)\), and each directed retention difference \(\ell_i=\Pr(L_i=1)-\Pr(I_i=1)\) for the supplied references is the simultaneous affine image of the common mode flow polytopes. The native isolated caps, containment bounds, directed retention bounds and global threshold tests can be checked as distinct claims on this one class, without requiring any factorization of local-event laws. The analogous simultaneous cap tests conditional on an independently nominated positive-mass prefix cylinder \(B\) use each numerator restricted to \(B\) and denominator \(\Pr(B)\). Conditional isolated reference risks in these comparisons are for the corresponding conditional mixtures. They require a separate premise to use an unconditional isolated cap as a conditional one.

The sharp bound for actual global risk is obtained by optimizing the total mass of the disjoint terminal patterns with \(G=1\) on one common flow. A failed universal bound has a single compatible witness law. Exact containment is equivalent to \(u=0\) throughout the feasible class. These conclusions require complete terminal scoring, including any local violations that can occur after a global event first fires. They make no claim of time- or space-efficient expansion from a succinct plant, discovery of trusted input constraints, unknown policy synthesis, or evaluation-to-deployment transport without source evidence.

\end{theorem}

\begin{proof}



\textbf{Source representation and policy legality.}

A physical prefix contains the information used to validate source transitions and conditional probabilities. A controller sees only the specified observation and its own permitted memory. In particular, the truth label attached to a leaf can be independent of the alarm without assuming probabilistic independence. For example, a source-checked receiving ledger can determine whether the actual delivery violated a governing route card while the alarm is computed from a pressure display. The semantic mode specifies which route card governs the scorer. Subsequent inference uses the ledger readbacks only where admitted and the truth label elsewhere only as a target. The monitor-alarm function and the source rule for truth remain distinct.

At a physical prefix \(h\), let \(a(e)\) be the callable action of a complete edge. For each action \(a\), the fixed-policy constraint is

\[
\sum_{e:\,a(e)=a}q_h(e)=\pi(a\mid O(h)).
\]

Here \(O(h)\) includes the authorized controller state as well as the admitted observation history. Identical controller inputs have identical action distributions, even when the physical state differs. Policy memory and any retained random seed required for this equation are part of the validated representation. The synthetic initial row uses the stated initial law or row cuts. At a positive prefix and an action of positive probability, the response kernel is obtained by dividing the edge weights in that action's block by the block total. The resulting probability law first makes the permitted policy choice, then executes the indicated complete source move with its full response. Zero-probability actions do not require speculative response kernels for this disintegration. Required operational support and legality on other prefixes are supplied by the source contract, including any alternatives needed by audited comparison policies.

D-35's enabled menu and touch check are separate from the mass constraints. A source permission check uses the ordinary precondition and the extensional authority state. A freeze excludes any touched coordinate of the given source atom over its supported successors. Uniform mode legality can be checked on the full declared candidate support before optimization, so an accepted action remains legal under every feasible allocation. A physical model with support-dependent menus may require additional modes or a sharper source support contract. Erasing a resource effect, equating two projected action labels, or adding scorer symbols to \(O\) would change the problem. D-35 source and hashes · Accepted source-conditioned T18.

\textbf{Scope reading after the blind intake.}

The set equalities and extrema apply to real law weights and fixed real affine statistics. The polynomial bit-complexity claim uses binary rational encodings for any coefficients, costs or cap values that the computation consumes, including the queried affine objectives and ratios. The input size includes those encodings.

Exact containment over the feasible class is equivalent to the absence of a terminal uncovered pattern with positive mass under any feasible law. Indeed, the uncovered probability is the sum of its terminal masses, all nonnegative. A leaf appearing in an explicit candidate tree may still have mass zero in every feasible law because of paired rows or affine evidence. A graph reachability check gives the same exact criterion if every syntactically reachable relevant terminal has positive mass under at least one compatible law. In the general flow model the test is maximization of uncovered mass.

\textbf{T25.J · Realization and sharp inference.}

\textbf{Exact flow representation.}

Fix a mode. For an admissible law on a tree, give each node its prefix probability and each edge its traversal probability. The children partition the parent event, yielding the flow equalities. At a positive prefix, multiplication of its conditional row constraints by the node probability yields the homogeneous inequalities. At a zero prefix all outgoing masses are zero, and every such homogeneous inequality holds. Admitted cross-context evidence holds by the completion's definition.

Conversely, fix a feasible rational-polytope point, possibly with real coordinates. At each positive node define \(q_h(e)=f_{h,e}/x_h\). Nonnegativity and flow conservation give a probability row and division gives every required row inequality. Start with mass one in each root and use these rows successively, separated into policy and response factors as above. At every depth induction gives its prefix masses exactly. At a positive parent the newly traversed edge has mass \(x_h q_h(e)=f_{h,e}\). At a zero parent the requested child masses are zero and the process never traverses its edges. Local rows at such null prefixes may even be empty. The finite law satisfies all verified cross-mass constraints and scores each realized path with the specified functions. Its law on complete paths is determined by the terminal masses. For multiple trees, the constructed tuple has one fixed mode and all cross-tree equations hold. A fully specified common-ticket execution can itself be a product tree when evidence concerns the pairing beyond the separate marginals.

Each coordinate lies in \([0,1]\) by root normalization and nonnegative conservation. Every mode feasible set is thus a compact polytope, including when it lies in a proper affine subspace. Taking precisely the realizable tuples over the listed modes yields the exact union. For any affine statistic \(a\), each feasible flow gives a law with its indicated value, and each compatible law contributes a feasible image point. This proves sharpness for the whole vector at once. Its representation as an affine image of a union suffices for linear queries and retains the possible gaps between fixed physical modes.

For example, to condition on a complete record class \(z\), sum leaf mass over the matching visibility predicate. A separate leaf audit tag can distinguish authenticated positives, authenticated negatives and unavailable truth. Label constraints are sums on the respective certified subsets. They can also refer to joint record and truth symbols or a legitimate calibration arm. A match between the source scorer and an authenticated label must be justified before treating that tag as truth about the intended population. The flow construction provides a witness for the conditional statistics only under that premise.

\textbf{Ratios, cross-context risk and certificates.}

Fix one feasible mode and write a point of its polytope as \(x=\sum_j\lambda_j v_j\) over its vertices. If \(d(x)>0\), affineness and \(n(v_j)=0\) at zero-denominator vertices give

\[
\frac{n(x)}{d(x)}
=\sum_{j:\,d(v_j)>0}\frac{\lambda_j d(v_j)}{d(x)}
\cdot\frac{n(v_j)}{d(v_j)}.
\]

The displayed summands are products of the nonnegative weights and the vertex ratios. Their weights sum to one. Thus the extreme defined values are attained at positive-denominator vertices. The segment between two extremizing vertices has positive denominator and its continuous ratio runs through the entire interval. A mode with maximum denominator zero contributes only undefined completions. Minimizing \(d\) on each feasible mode detects whether any denominator is zero. On the positive subset \(n/d\le c\) is equivalent to \(n-cd\le0\), and on the zero subset the latter is \(0\le0\). The proposed test therefore quantifies exactly over the defined values once a consistent class with a defined value has been established. The proof applies also to signed conditional event differences.

For explicit computation, write a mode as \(A x\le b\) with the bounds \(0\le x_i\le1\) included. Let \(n(x)=u^\top x+u_0\) and \(d(x)=v^\top x+v_0\). After checking that some \(d(x)>0\), solve the ordinary linear programs for the two extrema of \(u^\top y+u_0 s\) over

\[
A y\le bs,\qquad s\ge0,\qquad v^\top y+v_0s=1.
\]

Here \((y,s)=(x/d(x),1/d(x))\). At \(s=0\), the lifted coordinate bounds force every \(y_i=0\), contradicting the normalization. Hence every transformed feasible point has \(s>0\) and maps back via \(x=y/s\) to a positive-denominator completion. The transformed feasible set can be unbounded as \(d\) approaches zero. The preceding vertex argument gives attained finite ratio extrema.

For each feasible mode the maximum worst-context risk is the largest of the separate context maxima for that same mode. Maximizing over modes and the finite context index commutes. For the minimum, add one variable \(z\) and require \(r_j(x)\le z\) for every context \(j\) in the specified family, together with \(0\le z\le1\). Minimize \(z\) in each mode and take the smallest resulting optimum. Each optimum uses one common completion. The worst-context risk is continuous, so both outer extrema are attained when the feasible union is nonempty. The same conclusion holds for each scalar affine statistic. For the global threshold the upper endpoint determines universal passage, and the strict inequality of the lower endpoint with the threshold determines universal failure. Inconsistency is a separate output.

The feasibility and all the optimization problems have sizes polynomial in the explicit number of modes, nodes, edges, event coefficients and rational evidence rows. Their number is linear in the modes times the number of objectives plus context queries. Rational LP solves them in polynomial bit complexity. A rational optimum or feasibility witness can be checked directly against an explicit mode and reconstructed by the displayed row formula. Bounds on the coordinates of an optimal rational vertex follow by clearing denominators and bounding determinants of independent tight constraints. Upper-cap verification can also use a dual certificate for every feasible mode. With all rows in \(A x\le b\) and affine objective \(w^\top x+w_0\), any \(\lambda\ge0\) with \(A^\top\lambda=w\) and \(b^\top\lambda+w_0\le c\) certifies the cap. Strong LP duality supplies such a witness when the bound is valid. A certificate for an empty mode is \(\lambda\ge0\), \(A^\top\lambda=0\), \(b^\top\lambda<0\), from Farkas's lemma. These checks certify analysis of the encoded evidence. Deciding whether a row bound or a source menu belongs in the input requires its independent provenance. For fractional endpoints use the transformed LP and map the rational primal witness back to an allowed native law.

\textbf{T26.C · Paired continuation and optimal event coupling.}

\textbf{Exact family of couplings.}

Write \(h=(h_D,h_P)\). Given the native next-move probabilities \(K_j(h_j,e_j)\) for arm \(j\in\{D,P\}\), the product row equations at a positive paired prefix are

\[
\sum_{e:\,e_j\text{ is its }j\text{th move}}q_h(e)=K_j(h_j,e_j)
\quad\text{for every allowed native }e_j.
\]

If a particular native edge has no paired representative in the supplied support, this equality uses a zero sum. It can then force the paired prefix to be null. Independent paired-support or row-bound evidence is imposed on \(q_h\) in addition to these equations. All constraints are affine in the normalized row and have the stated homogeneous mass formulation.

Applying the flow realization gives a law on complete paired paths. Fix a positive paired prefix \(h\). For any next move in arm \(j\), the conditional probability equals \(K_j(h_j,e_j)\) by the equation. To extend to a fixed string of native moves, condition first on the possible paired next prefixes having positive mass and the specified projected move. By induction on remaining length, each such paired next prefix gives the same probability of the native suffix, namely the product of the native conditional probabilities along its prescribed native string. Multiplying the first factor and summing over the other arm's continuations proves the native conditional string probability starting at \(h\). A first factor of zero gives zero without conditional division. This argument applies at the root, yielding the stated complete native marginals, and applies starting at every positive paired prefix. Conversely, any coupling obeying the conditional relation integrates its rows to flows and satisfies the specified affine restrictions. An arbitrary coupling of the terminal events can have the same native marginals while failing this history-conditional relation.

A union \(B\) of disjoint paired-prefix cylinders of positive total mass partitions the conditioning event. On \(B\) the conditional law of each arm is the mixture of the just-proved native continuations with weights \(x_h/\Pr(B)\). Prefixes of mass zero contribute no term. For a variable horizon one may use the common semantic padding and carry terminated truth bits. Padding must leave the defined target untouched and preserve lawful stopping. All event probabilities on \(B\) are restricted sums of terminal masses.

\textbf{Backward recurrence and witness.}

For minimization give each terminal node its specified finite terminal cost. At a nonterminal node first mark every successor whose subtree is infeasible when given unit root mass. Impose zero mass on those successor edges in its normalized row polytope. If the restricted row is empty, mark the current node infeasible. Otherwise minimize \(\sum_e q_h(e)c_{\operatorname{child}(e)}\) using just feasible children. For maximization use the corresponding maxima. The comparison starts at the synthetic root.

Induct on distance from the leaves. Any admissible continuation at the node puts positive weight only on nonempty child continuation families, and conditional expectation over those children bounds its cost by the row optimum. Each restricted row polytope is compact. A row attaining its optimum followed by a chosen attaining child continuation on every positive edge gives a feasible law with value equal to the row optimum. If every row sends positive mass into an infeasible child then no valid continuation exists. This also proves the stated infeasibility recursion. Choose one attaining row per feasible positive node to reconstruct a paired certificate. Optional constraints tying distinct prefixes or populations can invalidate independent local choices, so the common global LP handles those cases.

There are at most two row LPs per nonterminal prefix for one extremum, counting an explicit feasibility check. For fixed rational leaf costs, the bit size of each finite child optimum is polynomial in the explicit subtree input. The whole subtree has the equivalent bounded flow LP with its original coefficients and a leaf-cost objective, and determinant bounds apply to that subtree optimum. This controls coefficient sizes in the backward recurrence. The running time is polynomial in the explicit paired tree, including the input describing the costs and restrictions. Histories can expand exponentially in an encoded horizon or number of coordinates. The result permits this expansion.

\textbf{Event fidelity, measurement and policies.}

At a terminal pair the indicator difference is \(1[V_D=1,V_P=0]-1[V_D=0,V_P=1]\). Taking expectations gives \(r_D-r_P=b^+-b^-\). The disagreement expectation equals \(b^++b^-\) and bounds the absolute difference. A minimum of zero is attained and yields a pathwise-agreeing coupling almost surely. Conversely such a coupling has objective zero. The minimizing cost is a bound from the existence of one relation preserving the prescribed native laws. The maximizing cost answers the different question of disagreement in every allowed physical pairing. For exact native laws the signed marginal risk difference is the same under each feasible coupling. A maximum discrepancy bound on a conditional cylinder across possible pairs or modes concerns those exact feasible pairs and is computed with the T25.J ratio test. The same binary identity holds after each positive conditioning.

Two individually legal policies can have different observations and take different callable actions. Native rows must describe those actual policies, including any allowed probe-dependent choices. A probe's truth-report channel can be scored separately, for example by the expectation of \(1[\mathrm{report}_P\ne V_P]\). This quantifies accuracy on the probe's own law. The discrepancy cost relates two named physical interventions and their specified policy laws. A faithful probe can have an uninformative report. A perfectly scored probe can change the actual risk. With a source-validated shared primitive, a pairing witness can record the provided physical correspondence. With unrestricted row couplings it witnesses a mathematical bound; a particular realized interventional pairing requires a separate source justification.

\textbf{T27.J · Exact conditional composition.}

Start from the full joint plant including the relevant physical shared resource, authority and communication state. A terminal label pattern \(s=(g,\ell_1,\ldots,\ell_k)\) specifies a disjoint block of terminals. Its mass is

\[
p_s=\sum_{h:\,\mathrm{pattern}(h)=s}x_h.
\]

These masses sum to one in each population tree. The actual global risk is \(\sum_{s:g=1}p_s\), the exact uncovered risk is \(p_{(1,0,\ldots,0)}\), and the local event masses are \(\sum_{s:\ell_i=1}p_s\). One joint flow retains all overlap and all constraints involving the joint trace, including construction through a memory of an earlier global event. A scorer that waits until the designated terminal correctly accounts for local events that fire later.

For an optional isolated reference \(I_i\), impose the fixed marginal row equalities for that coordinate at every positive paired prefix. The projection proof in T26.C applies to that coordinate with all remaining coordinates taken together. In particular the isolated terminal marginal is its known lawful native isolated risk. The uncertain actual joint coordinate can project from rows that vary with the cohistory. At each of its own positive prefixes, condition on the paired prefixes projecting there. The genuine projected conditional row is a convex average of the permitted projected rows at those paired prefixes. It remains in its native row polytope and keeps its prescribed policy marginal and source support. Thus there is an executable native joint law on physical histories with that conditional row. The paired description introduces no new controller observations. Additional native-prefix affine evidence, including tied population masses, is evaluated on the summed projected masses. An unknown shared parameter imposing non-affine conditional equalities would require a different completion representation or a proved reduction to listed modes and affine rows.

The retained risk difference \(\Pr(L_i=1)-\Pr(I_i=1)\) is an affine statistic of the same flow. The same applies separately to the reference risk, uncovered mass and global risk. Under the declared model the feasible joint image of all these numbers is exact by T25.J. To test a conjunction of upper claims, maximize each violation slack in every feasible mode. The conjunction holds if each such maximum is nonpositive. Otherwise one row-feasible completion gives a positive slack for a failed claim. The full vector image supports also verifying jointly attainable equality patterns and tradeoffs: impose the simultaneous affine measurements or threshold inequalities on one common flow and test feasibility. An upper containment tolerance certifies precisely the uncovered-event claim, so exact containment within the declared feasible class is equivalent to its maximum being zero on a consistent class.

For a nominated prefix cylinder \(B\), set \(d_B=\sum_{h\in B}x_h\) using disjoint prefix representatives. Every conditioned event numerator is the sum of its designated terminal masses extending one of those prefixes. A difference such as conditional local retention uses the difference of two such sums. These numerators vanish when \(d_B=0\), and the T25.J statement applies. In particular a universal conditional cap \(c\) uses the linear slack numerator minus \(c d_B\) at each mode. Conditioned reference risk uses the native continuation mixture already proved in T26.C. For example, when an isolated arm and actual plant share a fair ticket, a rare ticket can positively identify an isolated violation. The isolated unconditional small risk and its conditional probability one coexist. Reusing the unconditional number requires uniform source or evidence control over the relevant conditional continuation set.

Projection, support, independently attributed scoring, and composition evidence are separate obligations. If all information about labels collapses to the T27 abstract budgets, the baseline disjoint-atom realization attains its clipped sum. If rows, occupancies or common tickets constrain overlaps, the common flow can give strictly smaller bounds without changing the global-event rule or presuming independent local failures. An exact global upper bound can hold even when particular containment or local-retention cap tests fail. Concrete realizations and boundaries.

\end{proof}

\subsection{M06. Joint verdict and useful work}

\label{ext:m06}

Here the verdict and the work belong to one execution. A certificate of intrinsic possibility cannot replace a common callable command when the positive branch still owes a job.

\begin{theorem}[T29.J. Joint evidence, verdict and useful work]\label{ext:t29j}

An instance is a finite set of hidden states tagged by an immutable completion \(\theta\), a nonempty initial belief, and a finite public meter \(p\). A tag has a fixed binary target \(f(\theta)\). Public command names designate actual callable commands. For each \((p,s,a)\), a finite exact row gives its authorization, ordinary enablement and all positive supported successors. An authorized enabled command has nonempty support. An edge outputs a public record \(y\) containing its observation, acquisition charge \(c_y\) and work charge \(w_y\). Both charges are nonnegative and branch-specific; all members of a single record group have the same charges. The public update \(p'=F(p,a,y)\) is deterministic. A hidden successor retains its completion tag. Any future availability, price, relevant safety or utility effect is captured by these rows and states. A declared resource freeze removes each source-atomic occurrence with a changing supported successor before beliefs are formed. A command can reveal evidence, perform work or change future availability.

Mark safe states/edges and useful terminal states directly in the instance. An issuing policy acts on its permitted public history and must be safe on all its supported paths. A negative verdict is lawful exactly when every surviving tag has \(f=0\). A positive verdict is lawful exactly when all surviving tags have \(f=1\) and every surviving state is a useful terminal. One budget \(b\) caps actual accumulated work charges on every supported branch. For fixed move limit \(h\), minimize the maximum supported accumulated acquisition charge among valid policies, allowing \(+\infty\) for nonexistence. Invalid initial states and empty evidence have separate unsafe/inconsistent dispositions.

Let \(B_{a,y}\) be the union of exactly the genuine successors for the recorded branch. A common candidate command is executable at every state in \(B\), every supported edge is safe, and each branch has \(w_y\le b\). The optimum is zero at a valid stop, infinite if exhausted without a valid stop, and otherwise
\[
  V_h(B,p,b)
  =\min_a\max_{y:B_{a,y}\ne\varnothing}
       \{c_y+V_{h-1}(B_{a,y},F(p,a,y),b-w_y)\}.
\]
The recurrence gives an attaining supported deterministic policy when finite and a finite decreasing-horizon refutation of every failing acquisition threshold. Randomized choice gives no improvement in this supported worst-path model. A fixed public meter and fixed budgets allow the hereditary monotonicity under shrinking the nonempty belief. The decision problem with both finite integer budgets \(B_0,C_0\) and no stated depth cap also has a finite reachability solution. In the time-homogeneous rules on \((B,p,b,c)\), every proper success has a witness of depth at most \(|P|(2^{|S|}-1)(B_0+1)(C_0+1)\), including cases with zero-cost loops. Proper success means every supported play ends at a valid stop. Intrinsic target constancy alone need not provide a common positive policy.

The exact finite correspondence requires rows implementing an observation-consistent public command mapping over the entire compatible belief and preserving options, support, costs, actual safety/work labels and completion tags. A checked forward inclusion with possibly extra abstract rows yields a sound robust policy when concretely implementable. An exact minimum in the physical protocol additionally needs preservation of all its admitted alternatives and both-way supported coverage.

\textbf{Conditions and boundary cases.} Check nonemptiness and safety of the current belief before either verdict or another move at every node, including terminal children. Branch charges are public and successor information is complete. The residual-grid attractor uses integer charges or exact rational scaling. A safe edge can enter an unsafe terminal if edge and state markers differ. An unrestricted price continuum makes the stated residual integer grid invalid.

The callable command catalog is finite and exhaustive for the declared protocol. Legality and supported transitions use literal public command identities. Infinitely many one-move services of acquisition costs \(1+1/j\) give an unattained infimum.

An exact correspondence supplies observation-measurable translations of common commands and all actionable branches, both ways, with the same labels and charges. Coarsening the signal identifying which of two exclusive useful commands is executable can turn a finite optimum into infeasibility.

\end{theorem}

\begin{proof}

\textbf{Exact finite information state.}

A public history fixes \(p\), remaining horizon and work budget because the charges are in the record. For a concrete history in belief \(B\), an authorized selected command \(a\) can produce exactly its listed successors. An observed branch \(y\) has as posterior the union \(B_{a,y}\) of these successors. Each member has a predecessor in a genuine prefix followed by a supported row edge, so the union is exact. Conversely every actual extension lies in this union. Tags never switch. This also shows why checking an abstract action separately on each hidden state is insufficient unless the action corresponds to one literal callable public command. The freeze and all safety and authority checks are on original whole source-atomic occurrences before grouping by observation.

Test nonemptiness and initial safety first. A valid stop needs no more acquisition. A unanimous false belief certifies NO. A unanimous true belief yields YES only when every possible physical state has completed useful work. Both stops and future legality are functions of the current information state. A nonterminal unsafe possibility loses. Inputs with unsupported records or no compatible state do not give vacuous verdicts.

\textbf{Induction and exact optimizer.}

At depth zero, a valid stop attains cost zero, and all other states lose. For future depth, every deterministic policy first stops or chooses one command common to all current states. Reject it if it is absent, unauthorized, ordinarily disabled, unsafe on a supported edge, or if some supported record overspends work. For each remaining command and each nonempty record child, the fixed public charge \(c_y\) has already been incurred and the continuation sees exactly its child, remaining budget and updated meter. Every valid continuation costs at least the induction value from that child. The worst supported cost of choosing \(a\) is at least the maximum of those sums. Conversely choose a minimizing command and an induction-optimal child tree for each genuine record. The catalog is finite, hence the minimum is attained whenever finite. Concatenation is feasible and attains the stated recurrence. If every branch of a randomized policy is required to satisfy the supported constraints, retain one positively used command at each reached public history; this weakly reduces worst supported cost and preserves all other guarantees. No stochastic benefit for positive-risk objectives is asserted.

For an acquisition cap \(c\), a losing witness at \((B,p,b,h,c)\) first rules out both stops by showing an incompatible verdict or unfinished work, or gives unsafe or inconsistent input. For each public command it exhibits a prohibited state/edge or budget overspend, or one genuine record child that loses with remaining cap \(c-c_y\) and depth \(h-1\). The verifier recomputes the full child from the supported rows, checks the ranked horizon and covers every candidate command. Induction proves soundness and builds such a refutation whenever \(V_h>c\). A winning tree with cost \(v\) and losing threshold checks bound the optimum; the recurrence itself gives a checked table with exact costs, including infinity. At a smaller nonempty subset, a larger-cell winning tree remains lawful and its attainable observation children only shrink. Stop predicates remain valid, so the optimum weakly improves at fixed meter and budgets.

For finite integer budgets \(B_0,C_0\), incorporate both residuals into the finite belief graph, and require legitimate costs uniformly across every successor. Let \(W_0\) be the safe permitted stops. Add to \(W_j\) each safe state with an affordable common command all of whose true children lie in \(W_j\). A first-inclusion rank supplies a proper finite policy. On this finite alphabet, an observation tree in which every branch terminates has bounded depth. Otherwise some child at every level has unbounded remaining depth, yielding a nonterminating branch by recursive selection. Thus any proper policy's start is included. There are at most \(|P|(2^{|S|}-1)(B_0+1)(C_0+1)\) information states with both residuals. Stabilization within that many additions proves the depth bound. Zero-cost self-loops outside the attractor do not close. A grid bound needs actual integer charges, or exact rational scaling; the fixed-depth induction works with finitely many explicit real costs and exact comparisons.

\textbf{Transport, necessity and baseline advancement.}

For sound forward transfer, give an actual physical realization of every used abstract common command, and induct on the true run using observation-identifiable abstract histories and preservation of every actual outcome, charges, safety/utility facts and true tags. Extra abstract candidates can block a legitimate true decision, but any supported policy over the overapproximation transfers if all concrete labels and commands are faithful. For equality of optima, also transport physical policies back: the declared exact protocol accounts for every actual alternative, exposes the actionable records both ways and produces the same conditional branches and costs. Induction evaluates corresponding trees with the same outcome, costs and value. Merely tracking a physical observation into a lossy abstraction is enough for a restricted safe certificate when branches are covered, but gives no converse.

The baseline's T29.a solved a finite supported useful-goal task and T29.b bought pure external evidence before starting it. The joint objective here requires a truthful fixed target verdict and real utility on every positive branch and optimizes paid evidence inside the plant, with separately budgeted operational charges and state-changing information. The generic belief-state proof mechanism is a known adaptation. Three contracts for which command \(a_j\) is available except in contract \(j\) show that every pair can serve while the triple cannot. A hidden preliminary signal identifying which of two mutually exclusive commands is now legal separates a full meter from a lossy one. The model applies to a bound on supported outcomes and cannot certify a numerical chance risk without quantitative law information.

\end{proof}

\subsection{M07. A succinct service with a hidden key}

\label{ext:m07}

The finite description in this example names exponentially many possible service addresses. Each address query is cheap, but the allowed interface provides no bundled lookup of the installed exceptional key.

\begin{theorem}[T07.K. Succinct governed service separation]\label{ext:t07k}

Fix \(n\ge 1\), \(N=2^n\). An installation has an unobserved gateway freshness bit \(g\) and one unobserved exceptional service context \(k\in\{0,1\}^n\). Its succinct source description uses a fixed interpreter with \(n+1\) stored mode bits. An admitted public context is an arbitrary input \(x\in\{0,1\}^n\). Each useful service executes a gateway job and one context job after it. The source-atomic fast actions cost 1 and survive the R freeze exactly when respectively \(g=0\) and \(x\ne k\). The corresponding universally lawful slow actions cost 3 and 2. Each job advances in all supported outcomes; no other command skips or performs either job. The fixed acquisition target is the ability to perform both jobs in every admitted context under R freeze and work budget 3. Fast guards have an E/R source realization via independently specified representation blocks, and changing any exceptional key changes an actual frozen fast menu and optimum in a context. A universally successful schedule across the \(g=0\) cell is required as part of the positive construction.

The authorized acquisition interface accepts the unit-price gateway read or a unit-price point query for any \(x\), which returns \(\mathbf 1_{\{x=k\}}\). No other information about \(g,k\) is initially admitted, including from public authority code, which describes the conditional rules only. This acquisition comparison ends at certification, with no free trial executions or bulk queries. Full identity is the pair \((g,k)\). The same oracle and uniform wrong-or-abstain error requirement apply to both goals.

With exact answers, the target is precisely \(\mathbf 1_{\{g=0\}}\). Its minimum cost for error \(<1/2\) is one. Full-identity deterministic zero-error minimum worst price is \(N\). For a bound \(q\ge1\) on unit-price calls, the optimum randomized uniform identity success is \(\min(1,q/N)\), with optimum \(1/(2N)\) for \(q=0\). Bounds are in the \(n\)-bit succinct size and the implicit context namespace, including address/output capacity and the shared physical work model. For any randomized identity strategy with success at least \(1-\delta\) everywhere, finite worst expected acquisition \(t\) satisfies \(t\ge \ell(1-\delta-\ell/N)\) for every integer \(\ell\ge1\) by truncation; constant \(\delta<1/2\) forces \(\Omega(N)\).

In a separate shared noise regime, each lawful bit answer is independently flipped with the same known \(\eta\in(0,1/2)\) conditional on the queried bit and all preceding records. With uniform wrong-or-abstain risk at most \(\delta\in(0,1/2)\), a proper randomized or deterministic sequential audit identifies the target using worst expected at most \(a/\gamma\) calls, where \(r=(1-\eta)/\eta\), \(\gamma=1-2\eta\), \(a=\lceil\log(3/\delta)/\log r\rceil\). Full identity has an achievable worst expected bound \((Na+b)/\gamma\) for \(b=\lceil\log(3N/\delta)/\log r\rceil\). The noiseless expected lower bound still applies by simulation. Thus for fixed \(\delta,\eta\) identity needs \(\Theta(N)\) worst expected calls and the target has a constant bound. This is an expected-cost result, with a separate exact worst-path result in the noiseless regime. An explicit polynomial-total-cost deterministic complete pure identifying menu, executable as one legitimate sequence with bit capacity and fidelity as declared, upper-bounds full-identity cost by that sum. Installed-key disclosure, arbitrary trials or group tests are changes of interface.

\textbf{Conditions and boundary cases.} The one-call optimum is the fixed worst-path cost. For uniform error \(\delta<1/2\), the noiseless worst expected target cost is \(1-2\delta\) when a privately randomized read/guess mix is allowed. A call only on a randomized fraction of runs has expectation below one and retains the required error.

\end{theorem}

\begin{proof}

\textbf{Real governed differences.}

For each public \(n\)-bit context \(x\), the whole source-atomic fast gateway and service commands perform the real jobs and emit admitted progress. At the gateway keep a legitimate job-certificate pair in one equivalence block iff its guard bit \(g=1\). At a context stage use a context-indexed certificate pair in one block iff \(x=k\). Fast refreshes the designated pair by separating it if stale. Slow performs the same useful job with no representation refresh. Thus old source support and authority determine real R touches. Use open progress histories compatible with two opposite target worlds and distinct correctly labelled terminal histories, as in the accepted D-35 examples. Authority is a fixed declared predicate permitting precisely the commands of the service protocol under their ordinary preconditions. The R freeze removes an occurrence whose support includes a refresh. It does not delete progress or permit the remaining part of a mixed atom as an independently callable action.

The fixed interpreter compares \(x\) to the installed \(k\) using a length-\(n\) equality test, plus one Boolean gateway test. The parameter-bearing contract encoding is \(\Theta(n)\) bits. An \(n\)-bit public index describes each legitimate point-read command, and each answer is one bit. The interface contains a succinctly described complete point menu and no bundled or address-revealing key read. Reading a source rule conditionally on an unfilled slot does not reveal the installed value. A replacement \(k'\ne k\) changes apparent availability of fast in the real context \(x=k\) and changes its known-policy frozen context tariff from 2 to 1. Both costs and target versions are defined on the same physical service. If \(g=0\), one public policy chooses the fast gateway and slow context route for each \(x\), always completing at cost 3. If \(g=1\), even a fully informed policy pays at least \(3+1>3\). The service target depends exactly on \(g\), and every complete identity variable is operational.

\textbf{Exact deterministic and randomized hard-budget cost.}

Give all \(2N\) parameter pairs equal prior mass to bound uniform performance. At a deterministic decision-tree node before a gateway read, assume \(m\) feasible keys remain and \(q\) calls are allowed. With zero calls, one guessed pair can be correct. If \(g\) is known, \(q\) point reads can distinguish at most \(\min(m,q+1)\) keys: each positive answer singles out its key, and the all-negative branch permits one guess. If \(g\) is unknown and \(q\ge1\), at most \(2\min(m,q)\) pairs can be correct. A gateway query allocates \(q-1\) further calls to either bit and achieves the bound. A point query with both outcomes possible has one positive-key branch and a remaining-\(m-1\) branch, with \(q-1\) calls each. For \(q=1\), the two leaves give at most two successes. For \(q>1\), induction bounds the first by two and the second by \(2(q-1)\). The total number can never exceed \(2m\). An uninformative or repeated query wastes a call and cannot increase the bound. A randomized policy is a mixture of deterministic trees after fixing its coins, so its worst success is at most the same prior average \(q/N\) capped at one. For each fixed installation, reading \(g\), picking \(q-1\) distinct contexts uniformly, and uniformly guessing an untested key if all are negative reaches exactly this probability for \(1\le q\le N\). Reading all but one key after the gateway achieves sure identity at \(N\). With \(q=0\) the uniform guess among all \(2N\) matches the upper bound.

One gateway read gives target certainty; without a query a fair unknown bit forces worst error at least \(1/2\). For the expected-cost target optimum with bound \(\delta<1/2\), couple two opposite gateways with the same fixed \(k\) and same random tape before the first gateway query. Absent such a query, their outputs coincide. Under equal gateway prior the error probability is at least one half the probability of no gateway read, so worst expected cost is at least \(1-2\delta\). A gateway read with probability \(1-2\delta\) and otherwise a fair guess attains it. For full identity with uniform success at least \(1-\delta\) and worst expected cost \(t\), cut off after \(\ell\ge1\) calls, charging abstention as error. Success remains at least \(1-\delta-t/\ell\) by Markov's inequality. The hard bound gives \(\min(1,\ell/N)\), yielding the stated lower bound for every \(\ell\). Choosing \(\ell\) near \((1-\delta)N/2\) proves fixed-confidence order \(N\).

\textbf{Noisy replies under the same acquisition interface.}

For a queried bit let \(S_t\) be the sum of increments \(+1\) for reply 1 and \(-1\) for reply 0. Use stop at \(-a\) (declare 0) or \(b\) (declare 1), with positive integer barriers. Under bit 0, \(r^{S_t}\) is a unit-start nonnegative martingale because \(\eta r+(1-\eta)/r=1\). Under bit 1, \(r^{-S_t}\) is such a martingale. Bounded stopping or the nonnegative crossing inequality gives false acceptance probability at most \(r^{-b}\) and false rejection probability at most \(r^{-a}\). The stopped drift satisfies \(\mathbb E S_{\tau\land t}=-\gamma\mathbb E(\tau\land t)\) under 0 and \(+\gamma\mathbb E(\tau\land t)\) under 1, by summing predictable indicators of continued sampling. Since the stopped walk is in \([-a,b]\), monotone convergence proves finite expected time, at most \(a/\gamma\) for bit 0 and \(b/\gamma\) for bit 1. This proves proper stopping. For the gateway use equal barriers \(a\) and \(a\). For keys scan distinct contexts, applying \((a,b)\) each time until declaring a key, or abstaining if none. The uniformly valid conditional bounds apply after any permitted history. Union of the gateway error, the at-most-one tested true-key rejection, and at most \(N-1\) false accepts has probability at most \(\delta\) at the stated barriers. The worst expected total is at most \((Na+b)/\gamma\) including the gateway. The target-only strategy needs at most \(a/\gamma\) in expectation, with at most the same tolerated error.

An independent flip added to each exact oracle answer simulates the noisy interface for any adaptive algorithm. The exact lower bound therefore applies to noisy worst expected acquisition as well. The comparison is on fixed uniform confidence and a common price per legitimate call. Neither the martingale proof nor the mean sample bound supplies a deterministic noisy stopping deadline. Arbitrary dependence between flips would not justify the calculation.

\textbf{Natural interface boundary and comparators.}

If a polynomial-total-cost deterministic menu can lawfully be acquired as one sequence and its joint pure answers distinguish every completion, buying that menu identifies the full contract within its total price. On an explicit menu costing a polynomial amount no exponential optimal identity lower bound can follow in that model. Here the legal point menu has \(2^n+1\) implicit elements and there is no group query or remotely visible installed key. Allowing either changes the lower bound. An actual sampled service trial that reveals the key or a batch test is outside this pure acquisition comparison. It would need to share confidence, time and prices with the alternatives under a new model.

The point-query obstruction and biased random walk methodology have clear antecedents. The domain contribution is a succinct governed construction with real frozen service costs, one shared operational target and exact optimal identity comparisons under its declared interface. A stronger claim on arbitrary strategy-aware oversight or external novelty is unwarranted.

\end{proof}

\subsection{M08. Correlated services and actionable receipts}

\label{ext:m08}

An accepting path encodes the entire installation. Earlier jobs can reveal lawful choices later. Compressing past costs on a frontier works because the installed continuation depends on its graph boundary and the full public record.

An explicit layered directed multigraph has initial vertex \(v_0\), a declared set of terminal vertices at layer \(n\ge1\), and edges \(E_i\) from layer \(i-1\) to \(i\). A compatible installation is an entire accepting path. Remove edges outside all accepting paths. An edge \(e\in E_i\) specifies a mode with bit \(x(e)\in\{0,1\}\). Physically distinguishable alternatives may also differ in their allowed observations within stage \(i\). Modes can be arbitrarily correlated through the graph. At each stage one required job must be performed. The fast command costs integer \(a_i>0\) and survives the R freeze iff \(x(e)=0\). The common slow job costs integer \(a_i+w_i\) for \(w_i>0\) and is always allowed. Each actual job fully completes this stage and advances it along the installed path on its whole support. Neither a read nor any other permitted move completes or bypasses a job. Before and after each move there is a finite public meter \(p\), recording every observation-relevant carry affecting future rights and prices.

The mode-conditional source correspondence fixes the complete lawful public command names, both job tariffs, all positive supported observations, job events, local read availability and the resulting public-meter update. Any hidden branch consistent with the current \(e,p\) offers the same rows in these coordinates. A row is independently established from source guards, complete physical support and actual execution-event semantics for each possible mode. Those conditional facts do not disclose which mode is installed as an unpaid input. Commands are literal common callable acts in the declared physical protocol. The whole catalog of its admitted physical acts is covered. Reads have fixed finite rational nonnegative price per declared public meter and command, with prices publicly known. Reading may change \(p\) and can be locally enabled or forbidden based on \(e,p\). Job receipts also change \(p\), and may reveal facts about later jobs through correlations. Reads do not change the installed path, job guard, work tariff or stage. Job actions have zero acquisition charge, and read actions have zero work tariff. Supported outcomes of an action depend on the current \(e,p\); future path edges can affect the present only by their correlation with \(e\). Time is an explicit bound \(H\ge n\) on total moves. Work budget \(D\) is a nonnegative integer.

For path \(m\), the intrinsic decision target is \(f(m)=1\) iff its known-mode minimum support-robust work tariff for all jobs is at most \(D\). An actual audit/work policy must be lawful and finish within \(H,D\) on every branch. YES is lawful only on a nonempty cell of \(f=1\) paths after all jobs complete. NO is lawful only on a nonempty cell of \(f=0\) paths and does not require serving them. Value is the least worst supported acquisition charge for a correct policy, including \(+\infty\) for nonexistence. All branches of an issued result use concrete callable commands and real service.

\begin{theorem}[T30.CG.1. Exact frontier summaries]\label{ext:t30cg1}\label{ext:t30cg}\label{ext:t30g}

The intrinsic score of a path is \(\sum_i \kappa(e_i)\), with \(\kappa(e)=a_i+w_i x(e)\); capped sums use overflow \(D+1\). At a history at stage \(i\), store for each feasible current edge the attainable minimum and maximum capped intrinsic prefix score just before this edge, over all compatible prefixes that can be concatenated with it. Mark absent edges separately. The public meter, remaining work/move budgets and the edge-indexed extrema suffice exactly for legal common decisions, every supported observation, updates and the target verdict. Local read observations restrict current edges. A job observation restricts its edges, adds the respective \(\kappa(e)\) capped, and transfers extrema to the next edges by the shared boundary vertex. Initial extrema are zero on all possible first edges. Terminal extrema include the final edge. Precompute attainable suffix extrema on the DAG. The exact target minimum/maximum on a nonempty cell are found by joining current prefix extrema, current edge score and path suffix extrema. Extremizing paths can be recovered using actual compatible representatives. The abstract state never requires filling in intermediate unattained sums.

\textbf{Conditions and boundary cases.} Every allowed operation is a modeled job or stage-local nonservice move with a full callable row, public quote, meter effects and supported records. Nonservice moves cannot change the installed path or serve jobs. Each branch key includes the full actionable next public record, including the next meter when a visible letter alone does not determine it. Use fixed graph-compatible initialization, ordinary fast enablement on bit-zero modes, and per-installation or full-positive prior conditions for zero-risk equivalence. Additional hidden operations or actionable observations invalidate minimum or failure claims. A mixture can hide a failing possible installation if its prior weight is zero.

The eight simultaneous obligations use the exact joint online deadline objective. T30.C is an exact static pure-target objective with total callable independent readbacks and its own timing. Static target certainty can leave no common useful policy on an intrinsically positive cell. The original T30.S predeployment menu can distinguish differently hidden later bits at the same first edge. A stage-local deterministic model reproducing those full first-stage read signatures may require a separate first edge for each distinct signature.

The initial public meter is fixed. Initial evidence must be incorporated into the exact accepting path class. A known restriction on entire first-edge cylinders can be handled by the admitted first edge set. An initial signal cutting a later-edge cylinder can make a feasible compressed continuation spurious unless the graph already represents that restriction.

Fast is ordinarily enabled and lawful at every active meter exactly on bit-zero modes after the freeze. The slow route is callable at every active meter. Every job advances and completes real service on its entire positive support. Every state-local callable action catalog is finite. Surviving a freeze alone cannot establish ordinary enablement. An undisclosed disabled fast route changes the intrinsic score.

Sure correctness equals zero violation probability for each compatible installed path under conditional laws with all listed edges positive. A mixture interpretation requires positive initial mass on every compatible path. A prior concentrated on a working installation can hide a supported failure confined to a different compatible installation.

\end{theorem}

\begin{theorem}[T30.CG.2. Joint optimal policy and refutation]\label{ext:t30cg2}

A finite budgeted minimax recursion over these compressed public states gives the exact least worst acquisition charge for the physical serial protocol. A finite value produces a real correct and useful policy on YES paths. Infinite values or values above a given cap give checkable refutations with coverage for every possible common first move. If the score cell is mixed, retain explicit feasible opposite-target paths. A verified losing joint-policy refutation has a separate interpretation from inability to identify a pure-read target. With exact guards and support, for every coherent finite stochastic law positive on each listed conditional supported transition, sure success is equivalent to zero probability of violation. The result makes no numerical inference about nonzero chance risk from support alone.

Writing \(W=\max_i |E_i|\), \(P=\sum_i |P_i|\) including terminal meters, and \(K=1+(D+2)(D+3)/2\), there are at most \((H+1)(D+1)P K^{\max(1,W)}\) memo states (terminal vertices can be folded to an aggregate terminal summary). Per-state work is polynomial in the explicit graph, supported row tables and integer or rational bit lengths. Reported state-count dependence is pseudo-polynomial for fixed width and numeric \(D,H\). The positive claim compares to an explicit path-catalog state or a full subset of all attainable \((e,\text{prefix-score})\) pairs. It gives no unconditional domination of other symbolic algorithms and does not include physical anchor-establishment cost.

\end{theorem}

\begin{theorem}[T30.C.3. Correlated pure static disagreement oracle]\label{ext:t30c3}\label{ext:t30c}

In a separately declared pure predeployment read interface, each eligible query has a deterministic answer from its addressed edge and a fixed additive price. A set of queries resolves the intrinsic threshold target on all feasible answer cells iff the product of two copies of the path graph, constrained to equal selected read answers at each layer and carrying both capped prefix scores, has no terminal state with opposite threshold bits. A surviving paired state yields two actual opposed accepting paths. For fixed reads the operation count is polynomial in \(\sum_i |E_i|^2\), local query tables and \((D+2)^2\). Exact implicit weighted hitting-set generation gives the minimum static target audit or an equal-full-menu obstruction, with the outer optimization cost separate. Static identification does not imply a shared budget-feasible useful policy. The local-coordinate hybrid rule of the Cartesian fragment is invalid on general accepting paths.

\textbf{Conditions and boundary cases.} Each static query is total and callable for every initial candidate installation in its separate predeployment interface. It has a pure deterministic response independent of other selected queries. A mode-dependent request denial either is an admissible observable outcome with a priced test, or excludes the request from this total menu.

\end{theorem}

\begin{proof}

\textbf{I. Intrinsic closure and contextual correspondence.}

Let \(m\) be any entire allowed path. Physical correspondence proves for each mode that fast is lawful, ordinarily enabled and succeeds with tariff \(a_i\) iff the guard survives the declared freeze, and that slow is always a common callable job at tariff \(a_i+w_i\). No read changes either condition, supplies service, or permits bypass. Every permissible success therefore pays at least \(\kappa(e_i)\) at each required stage. Knowing \(m\) one can choose the less costly available legal command at every stage, and reach the actual terminal service event on its complete support in \(n\le H\) moves. This proves the exact message-independent intrinsic cost and its threshold target. The conditional anchor states primitive source guards, genuine action availability, supported outcomes, charges, phase progress and actual events. It is checked over each possible mode with no initial read of which mode is installed.

Starting from a nonempty exact path class and the declared initial public meter, induct on actually observed prefixes. Once a candidate concrete path reaches the current stage and public meter, its entire concrete conditional positive observation and successor-image support equals the declared row. Each structured observable continuation is actual for that mode by appending some concrete supported step to an existing prefix. Every real branch maps the other way because the whole allowed physical catalog and support are covered. The meter preserves all actionable records. Choosing a command once for the full information state is implementable because the names are literal, fully authorized calls, with all supported transitions certified. A physical proof of totality and faithful linkage to actual work is required in addition to computing with the tables. Forward containment suffices for robust certificates if the data overapproximate actual support and calls remain sound. The exact optimum, negative witnesses and tight source-existence statements require exact rows and a realizable path class.

\textbf{II. Frontier invariant and exact target.}

Consider one allowed public history at stage \(i\). Before taking the current edge, prior completed prefixes have produced their actual measurements and receipts. Each path in the residual class contains one \(e\in E_i\). For a fixed such \(e\), membership of a prefix is determined by its completed edges, current edge and the past observations. Once the prefix and \(e\) are feasible, every accepting suffix following \(e\) is feasible under that same history. The graph allows all these concatenations and the probabilities and side fields controlled by the supported model have imposed no condition on the as-yet unobserved suffix. This holds initially because the admitted graph includes all actual initial evidence. It propagates because every read and job receipt in the model depends on the current edge and public meter, with its whole support certified at every compatible prefix.

Let \(C(t)=\min(D+1,t)\). For each feasible \(e\), store the least and greatest attained \(C(t)\) for the intrinsic prefix sum \(t\), and use a distinguished absent marker otherwise. When a read emits \(y\), keep exactly those entries whose mode-conditioned row includes \(y\). All prefix scores of a kept entry remain feasible because \(y\)'s support depends only on \(e,p\). When a job emits \(y\), remove the excluded entries, compute \(C(l+\kappa(e))\) and \(C(u+\kappa(e))\) for each kept entry \([l,u]\), then merge these attained extrema at each destination. For a next edge starting there, copy the resulting pair. The capped sum identity \(C(C(t)+k)=C(t+k)\) for nonnegative \(k\) and monotonicity ensure both exact endpoints of every image. Union extrema are attained in one old contributing entry. The induction neither inserts intermediate prefix scores nor drops a jointly feasible endpoint. At the terminal aggregate final extrema across the reachable final vertices using the same rule.

Compute \(\sigma_{\min}(v),\sigma_{\max}(v)\) for all accepted graph suffixes by backward min/max of capped \(\kappa\)-increments, using zero at accepting final vertices. The target lower endpoint at stage \(i\) is the minimum over present \(e\) of \(C(l_e+\kappa(e)+\sigma_{\min}(\operatorname{head}(e)))\), with the matching max using upper endpoints. By suffix compatibility each minimizing or maximizing prefix can be concatenated with that optimizing suffix. Representative prefixes and suffixes can be replayed with backpointers to obtain actual path and record witnesses, including opposed paths if the cell is mixed. On replay take the history that reached the chosen compressed state; coordinate extrema need not share the same prefix across different edges.

\textbf{III. Exact joint audit/work value.}

A lawful next selection, its real work cost, acquisition price and the entire set of possible public outcomes are functions of the supported current edges and \(p\). The next compressed state has just been shown exact for every actual history realizing it. Stop eligibility is a function of the endpoint comparisons and actual stage progress. Carry residual time and work budget in the state. Unsafe/unaffordable actions are eliminated using every supported record. Stop with NO iff the lower score exceeds \(D\), or YES iff all true paths meet the deadline and every job is done. For another move, take the minimum over common callable commands of the maximum over genuine receipts of quote plus next value. Horizon decreases. The same first-command/branch induction as in T29.J proves minimax exactness even when different previous transcripts lead to the same compact state. At winning states one can pick minimizing commands and achieve useful service on every positive path. At a failing threshold every feasible move has a supported losing successor, giving a finite failing certificate with exact branch coverage. Compute values and any counterexample along an actual possible history, so the implementation cannot choose different paid histories based on hidden state. A mixed cell produces opposed exact targets, while a homogeneous positive cell can still lose from lack of a common policy.

There are \(D+2\) capped scores and \(K=1+(D+2)(D+3)/2\) possible absent/ordered endpoint entries for each edge. With edge-width \(W\), \(P=\sum_i |P_i|\), \(H+1\) times and \(D+1\) work values, the stated storage bound follows. Every state update examines the declared local rows and successor edges and takes extrema; it is polynomial in those explicit sizes. Known-mode scores and work budgets are integer. Exact rational prices have bit length bounded by repeated arithmetic on up to \(H\) table charges. For fixed width and polynomial numeric time and deadline, this is a polynomial-state algorithm even on DAGs with exponentially many paths. The full possible score-set state could have \(2^{W(D+2)}\) different subsets before other dimensions; reachable exact-state counts can be lower, and decision-diagram and antichain variants may compress further. The proof gives a structural bound relative to the supplied graph and verified finite tables. Work to derive the graph, establish exact physical support or independently verify the anchor is a separate resource. Binary score hardness from T10 remains.

Support-robust failure is equivalent to positive failure probability for an installed path with any coherent full-positive conditional law on its finite supported tree: a finite failing branch has the product of positive conditional probabilities. Conversely no failing support implies no failure. Universal quantification is over each compatible installation. Integrating over an initially full-positive prior preserves the same equivalence, while a prior that puts no mass on an incompatible-to-success path does not.

\textbf{IV. T30.C · correlated static oracle.}

In the total deterministic pure predeployment menu, fix a set \(Q\) of selected stage-local queries. Form the product graph of two accepted paths and discard each edge pair whose \(Q\)-answers differ in its layer. Tabulate reachable pairs of boundary vertices and both capped accumulated scores. The roots have pair \((0,0)\); each surviving paired edge advances both endpoints. Induction proves every retained state comes from a real compatible pair with equal acquired answers, and every such pair is retained. Looking for one terminal side at most \(D\) and the other in overflow gives exactly the opposing pair condition. Keep actual predecessors for witnesses. With two bins of size \(D+2\) the cost is bounded by a polynomial combining \((D+2)^2\sum_i |E_i|^2\) state-edge updates and local answer tests. There is no need to materialize all models. Different global paths can lack single-coordinate hybrids, so the T30.S local-crossing proof cannot simply be reused.

For fixed-price static optimality, begin with no pair constraints. Minimize exact weighted cover of the disagreement sets seen so far and use this oracle to demand an opposed equal-answer pair. It either supplies a fresh unhit edge, or verifies sufficiency and the current minimum is a lower bound for any sufficient plan and is attained by the current plan. An empty edge witnesses full-menu inaccessibility. The finite query list bounds the number of strict refinements, but not polynomially in its size in general. This is the accepted implicit-cover logic with a new structured service oracle. A purchased target certificate alone supplies no common useful policy. The joint optimum of section III verifies that separate obligation in the same underlying service model with its own authorized timing.

\textbf{V. Concrete strengthening and limits.}

Use three jobs of base tariff 1 and extra tariff 1. The first job is fresh in both installations, with an auxiliary mode tag \(j\in\{0,1\}\). For \(j=0\) the last two freshness bits are \((0,1)\), and for \(j=1\) they are \((1,0)\). The first common fast job reports the independently declared tag during actual useful service; later receipts can be uninformative. With \(D=4,H=3\), both paths have intrinsic score 4. The public controller performs the first fast job and uses the receipt to choose the correct two remaining commands, taking zero paid reads and serving all jobs at cost 4. If the first receipt is constant and no informative reads exist, the second job must be slow, reveals nothing and leaves the third slow. Budget 4 is then infeasible although all paths are intrinsically positive. This is a genuine correlated work-and-information effect. The later fast guards change physical frozen options in two real jobs. The Cartesian replacement includes both stale modes simultaneously, changing the target.

For a mixed cell, take two jobs with base tariff 1 and slow premium 1. The first job is always fresh and emits a legitimate meter enabling a paid second-stage query, with no earlier queries. The second job's bit is unknown. \(D=2,H=3\). Without that query every transcript before a second-stage command is identical, so neither accurate NO nor a common affordable service is available. Paying the query, then doing fast for bit zero or issuing NO for bit one, realizes exact worst acquisition price equal to the quoted price. Both examples require an authorized receipt and keep the installed guard hidden until a lawful observation reveals it.

The proof considers threshold correctness, explicit witnesses, completeness, exact physical correspondence, optimum worst total read cost, conditional physical transfer, real service and a bounded representation simultaneously under the same online serial class. It broadens T30.S to correlations and useful feedback, changing the audit timing and objective explicitly. The fixed predeployment static result is T30.C separately. This support-only argument carries no positive-risk bound for an unmeasured side event.

\textbf{VI. Positive-risk and strategic scope obstruction.}

Consider an actual useful one-stage service always callable at unit cost, with a terminal side event \(V\) whose actual occurrence is independently well defined. Two permitted installed laws have \(P(V=1)=1/4\) and \(3/4\). All side-event alternatives are positive at both installations. The permitted evaluator and any admitted joint query/work policy record exactly the same commands, outcomes, prices and progress, omitting the side-event. It follows by induction on public histories that all such transcripts have the same law, including randomized decisions and common stopping, despite opposite risk labels at cutoff \(1/2\). Summing the two wrong-or-abstain probabilities gives at least one. Source support equality proves neither chance-regime membership nor useful-event risk beyond zero in this setting. This obstruction identifies the missing premises. For example, Theorem~\ref{adv:ug3} supplies both a checked physical risk mechanism and an informative lawful observation on a particular service. Its conclusion does not transfer to the equal-law pair here. The impossibility includes informative physical operations only when their actual observed joint rows still coincide. It cannot be inferred from equal benchmark marginals alone.

\end{proof}
